\documentclass[a4paper,12pt]{report}

\usepackage[T2A]{fontenc}
\usepackage[utf8]{inputenc}
\usepackage[english,russian]{babel}
\usepackage{lmodern}
\DeclareFontFamily{T2A}{lmr}{}
\DeclareFontShape{T2A}{lmr}{m}{n}{<->ssub*cmr/m/n}{}
\DeclareFontShape{T2A}{lmr}{bx}{n}{<->ssub*cmr/bx/n}{}
\DeclareFontShape{T2A}{lmr}{b}{n}{<->ssub*cmr/bx/n}{}
\DeclareFontShape{T2A}{lmr}{m}{it}{<->ssub*cmr/m/it}{}
\DeclareFontShape{T2A}{lmr}{m}{sl}{<->ssub*cmr/m/sl}{}
\usepackage{graphicx}
\usepackage{setspace}
\usepackage{fancyvrb}
\usepackage{appendix}
\usepackage{listings}
\usepackage{color}
\usepackage{amssymb}
\usepackage{amsmath}
\usepackage{hyphenat}
\usepackage{subcaption}
\usepackage{hyperref}
\usepackage{float}
\usepackage{booktabs}
\usepackage{array}
\makeatletter
\renewcommand\@biblabel[1]{#1.}
\makeatother

\newcommand{\myrule}[1]{\rule{#1}{0.4pt}}
\newcommand{\sign}[2][~]{{\small\myrule{#2}\\[-0.7em]\makebox[#2]{\it #1}}}

% Поля
\usepackage[top=20mm, left=30mm, right=15mm, bottom=20mm, includefoot]{geometry}
\usepackage{indentfirst}
\footskip 20pt

\definecolor{gray}{rgb}{0.5,0.5,0.5}
\setcounter{secnumdepth}{4}

\bibliographystyle{gost780s}  %% стилевой файл для оформления по ГОСТу

\lstset{
	basicstyle=\footnotesize,           % the size of the fonts that are used for the code
	numbers=left,                   % where to put the line-numbers
	numberstyle=\tiny\color{gray},
	captionpos=b,                   % sets the caption-position to bottom
	breaklines=true,                % sets automatic line breaking
}


\begin{document}

	%%%%%%%%%%%%%%%%%%%%%%%%%%%%%%%
	%%%                         %%%
	%%% Начало титульного листа %%%

	\thispagestyle{empty}
	{ \large
		%%% Название работы %%%
		\begin{center}
			\renewcommand{\baselinestretch}{1.2}
			\Large
			{\Large Анализ зашумленных и частично наблюдаемых данных судовых навигационных систем}
		\end{center}
		\medskip%

		\vspace{3\bigskipamount}%
		\begin{flushright}
			\begin{minipage}{0.5\textwidth}
				\begin{flushleft}
					Гашков С.П.
				\end{flushleft}
			\end{minipage}
		\end{flushright}

		\vfill%
	}

	%%% Конец титульного листа  %%%
	%%%                         %%%
	%%%%%%%%%%%%%%%%%%%%%%%%%%%%%%%

	% Размер шрифта (14пт)
	\large

	% Межстрочный интервал (полуторный)
	\onehalfspacing

	\clearpage
	\renewcommand{\thesection}{\arabic{section}}
	\graphicspath{{../pict/}}

	\section{Постановка задачи}
	Дан временной ряд
	$\left\{\left(\mathit{time}_i,\, \mathit{lat}_i,\, \mathit{lon}_i,\, \mathit{sat}_i \right) \right\}_{i=0}^{n}$, где:\\
	$time_i: \; time_{i+1} - time_i = \Delta t = const \qquad \forall i \in [0..n-1]$ --- дата и время, в которое произошло очередное измерение (по умолчанию $\Delta t = 10$~с),\\
	$lat_i: \; lat_i \in [-90.0 \; \ldots \; 90.0]$ --- широта, полученная датчиком GPS в момент времени $time_i$,\\
	$lon_i: \; lon_i \in [-180.0 \; \ldots \; 180.0]$ --- долгота, полученная датчиком GPS в момент времени $time_i$,\\
	$sat_i: \; sat_i \in \{0 \; \ldots \; N\}$ --- количество спутников, видимое датчиком GPS в момент времени $time_i$,\\
	являющийся записью геопространственной телеметрии судна.

	Основная задача: составить временной ряд меток $\left\{y_i\right\}_{i=0}^{n}$, где:\\
	$y_i =
	\begin{cases}
		\mathrm{move},   & \text{$i$ --- точка со штатным движением судна}\\
		\mathrm{anomaly}, & \text{$i$ --- точка с аномальным движением судна (помехами GPS)}
	\end{cases}$,
	являющийся результатом классификации точек на штатные и аномальные.

	В целях исследования, обработки, обучения и тестирования моделей заказчиком предоставлены 3 временных ряда (реальных телеметрий судов), соответствующих 3 разным судам (разных типов, конфигураций и районах движения) за период 3 месяцев следующего вида:
	$\left\{\left(\mathit{time}_i,\, \mathit{lat}_i,\, \mathit{lon}_i,\, \mathit{sat}_i,\, y_i\right) \right\}_{i=0}^{n}$, где $y_i$ является эталонной ручной разметкой, произведённой специалистом предметной области со стороны заказчика.

	\section{Список метрик}
	Основная цель заказчика --- максимизация достоверности статистики движения судна (в частности, пройденного пути) и положения судна в конкретный момент времени. В связи с указанной целью, а также тем, что данная задача является задачей классификации, в работе используется следующий набор метрик (в порядке приоритетности):
	\begin{enumerate}
		\item Геопространственные метрики,
		\item Временные метрики,
		\item Общие метрики классификации.
	\end{enumerate}

	\subsection*{Ограничение}
	В рамках поставленной задачи истинные точки траектории движения судна известны только для точек $y_i \ne \mathrm{anomaly}$. В точках $y_i = \mathrm{anomaly}$ истинные точки траектории неизвестны. Ввиду этого вводится предположение о том, что в аномальных точках происходит равномерное, прямолинейное движение по прямой, соединяющей две ближайших не аномальные точки. Предполагается, что такой способ заполнения данных, примененный и к предсказаниям модели, и к эталонным данным, обеспечит корректный расчет геопространственных метрик. В результате введения данного ограничения формируется следующий порядок расчета геопространственных метрик:
	\begin{enumerate}
		\item Для эталонного ряда строится траектория с координатами на всех индексах $i \in [0..n]$:
			\begin{itemize}
				\item В точках, где $y_i \neq \mathrm{anomaly}$ значения $\mathit{lat}_i$, $\mathit{lon}_i$ сохраняются из исходных данных;
				\item В точках, где $y_i = \mathrm{anomaly}$ координаты заполняются линейной интерполяцией между ближайшими не аномальными соседями. Для каждого такого $i$:\\
				$j: j < i,\; j=\max,\; y_j \neq \mathrm{anomaly}$ --- последняя не аномальная точка до $i$;\\
				$k: k > i,\; k=\min,\; y_k \neq \mathrm{anomaly}$ --- первая не аномальная точка после $i$;\\
				$\alpha_i = \dfrac{time_i - time_j}{time_k - time_j}$ --- доля пройденного времени от $time_j$ до $time_k$;\\
				$\mathit{lat}_i = \mathit{lat}_j + \alpha_i(\mathit{lat}_k - \mathit{lat}_j)$, \quad $\mathit{lon}_i = \mathit{lon}_j + \alpha_i(\mathit{lon}_k - \mathit{lon}_j)$.
			\end{itemize}
		\item Аналогично рассчитываются точки для трека, обработанного моделью (в качестве $y_i$ используются предсказания модели).
		\item Рассчитываются геопространственные метрики.
	\end{enumerate}

	\subsection*{Геопространственные метрики}
	Для расчета геопространственных метрик используется геодезическое расстояние между точками траектории в один и тот же момент $time_i$:\\
	$d_i = d_{\mathrm{geo}}\bigl((lat_{truth,i},\, lon_{truth,i}),\, (lat_{prediction,i},\, lon_{prediction,i})\bigr)$, где:\\
	$lat_{truth,i}$, $lon_{truth,i}$ --- координаты из эталонных данных,\\
	$lat_{prediction,i}$, $lon_{prediction,i}$ --- координаты из предсказания модели,\\
	$d_{\mathrm{geo}}$ --- геодезическое расстояние.\\
	Сами геопространственные метрики представлены следующим набором метрик:
	\begin{enumerate}
		\item $RMSE$\\
		Среднеквадратичное отклонение (в метрах) между соответствующими точками траекторий только в точках с ненулевым отклонением (точки с $d_i = 0$ в сумму не входят, чтобы не занижать ошибку там, где траектории совпадают).\\
		Обозначим $I=\{i : d_i > 0\}$, $|I|$ --- число таких индексов. Тогда
		$RMSE=\sqrt{\frac{1}{|I|}\sum_{i\in I} d_i^2}$, где:\\
		$d_i$ --- геодезическое отклонение между $i$-ми точками траекторий (в метрах);
		Если $I=\varnothing$, метрика не определена.\\
		\item $Max(d_i)$\\
		Максимальное отклонение (в метрах) между точками траекторий.
		\item Distance loss\\
		Разница суммарных длин траекторий.
		$DL=|L_{\mathrm{truth}}-L_{\mathrm{prediction}}|,  \qquad L(T)=\sum d_{\mathrm{geo}}(p_i,p_{i+1})$, где:\\
		$p_i$ --- $i$-я точка траектории с координатами $(lat_i, lon_i)$;\\
		$d_{\mathrm{geo}}(p_i,p_{i+1})$ --- геодезическое расстояние между соседними точками (в метрах);
	\end{enumerate}

	\subsection*{Временные метрики}
	\begin{enumerate}
		\item Detection Delay\\
		Задержка от начала истинного эпизода аномалии до первой обнаруженной аномалии внутри него. Под \textbf{обнаружением} аномалии в точке $i$ понимается предсказание модели $\hat{y}_i = \mathrm{anomaly}$.\\
		Для непрерывного интервала $S_k=\{i : y_i = \mathrm{anomaly}\}$:\\
		$\mathrm{Delay}_k=time_k^{\mathrm{detect}}-time_k^{\mathrm{start}}$, где:\\
		$time_k^{\mathrm{start}}=\min_{i\in S_k} time_i$ --- момент начала аномалии,\\
		$time_k^{\mathrm{detect}}=\min\{time_i : i\in S_k,\ \hat{y}_i = \mathrm{anomaly}\}$ --- момент первого обнаружения аномалии на участке.\\
		Если внутри $S_k$ аномалия не обнаружена ($\hat{y}_i \neq \mathrm{anomaly}$ на всём участке), событие считается \textbf{не обнаруженным}; $\mathrm{Delay}_k$ предполагается равным длине эпизода $time_k^{\mathrm{end}}-time_k^{\mathrm{start}}$, где $time_k^{\mathrm{end}}=\max_{i\in S_k} time_i$.\\
		Агрегаты по задержке (среднее, медиана, 95-й перцентиль, максимум) считаются по всем эпизодам с $\mathrm{Delay}_k > 0$: по опозданиям (обнаружение есть, но не в первой точке $time_k^{\mathrm{start}}$) и по не обнаруженным эпизодам (штраф --- длина $S_k$). Эпизоды с обнаружением в первой точке ($\mathrm{Delay}_k = 0$) в эти агрегаты \textbf{не входят}, чтобы нули не занижали среднюю задержку.\\
		Дополнительно фиксируются счетчики по всем истинным эпизодам $S_k$:\\
		$N_{\mathrm{on\,time}}$ --- число эпизодов с обнаружением в первой точке ($\mathrm{Delay}_k = 0$);\\
		$N_{\mathrm{late\,or\,missed}}$ --- число эпизодов с опозданием ($\mathrm{Delay}_k > 0$ при детекции) или без детекции внутри $S_k$.\\
		Если нет ни одного эпизода с $\mathrm{Delay}_k > 0$ (только своевременные детекции или нет истинных эпизодов), агрегаты задержки не определены.
	\end{enumerate}

	\subsection*{Общие метрики классификации}
	Классические метрики задачи классификации. Положительный класс --- $\mathrm{anomaly}$.\\
	$TP=\sum_{i=0}^{n}\mathbf{1}[y_i=\mathrm{anomaly}\land\hat{y}_i=\mathrm{anomaly}]$, \\
	$FP=\sum_{i=0}^{n}\mathbf{1}[y_i=\mathrm{move}\land\hat{y}_i=\mathrm{anomaly}]$, \\
	$FN=\sum_{i=0}^{n}\mathbf{1}[y_i=\mathrm{anomaly}\land\hat{y}_i=\mathrm{move}]$, где:\\
	$y_i$ --- эталонная метка, \\
	$\hat{y}_i$ --- предсказание модели.
	\begin{enumerate}
		\item Precision\\
		Доля предсказанных аномалий, которые действительно являются аномалиями.\\
		$\mathrm{Precision}=\dfrac{TP}{TP+FP}$.
		\item Recall\\
		Доля истинных аномалий, которые модель обнаружила.\\
		$\mathrm{Recall}=\dfrac{TP}{TP+FN}$.
		\item $F_1$\\
		Гармоническое среднее точности и полноты.\\
		$F_1=\dfrac{2TP}{2TP+FP+FN}$.
		\item MCC\\
		Коэффициент корреляции Мэттьюса по точкам с оценкой: учитывает все четыре ячейки таблицы сопряжённости, определение --- в разделе «Описание метрик». Используется при выборе порога и показывается в итоговых таблицах.
	\end{enumerate}


	\section{Предварительная обработка данных}
	В рамках предварительной обработки данных удаляются заведомо ошибочные записи, а также добавляются дополнительные поля.

	\subsection*{Удаление записей}
	Из временного ряда удаляются следующие записи:
	\begin{enumerate}
		\item $\left\{\left(\mathit{time}_i,\, \mathit{lat}_i,\, \mathit{lon}_i,\, \mathit{sat}_i\right) \right\}_{i=0, i \in {lat_i = Null \lor lon_i = \mathrm{Null}}}^{n}$,
		\item $\left\{\left(\mathit{time}_i,\, \mathit{lat}_i,\, \mathit{lon}_i,\, \mathit{sat}_i\right) \right\}_{i=0, i \in {lat_i = 0.0 \land lon_i = 0.0}}^{n}$ --- по причине того, что некоторые модели GPS-датчиков вместо заведомо ошибочных координат (например, нарушающих ограничения $lat_i \in [-90.0 \; \ldots \; 90.0]$, $lon_i \in [-180.0 \; \ldots \; 180.0]$) записывают координаты (0.0, 0.0).
	\end{enumerate}

	\subsection*{Добавление полей}
	Во временной ряд добавляется поле $in\_water$:\\
	$in\_water  =
	\begin{cases}
		\mathrm{true},   & \text{$i$ --- точка, находящаяся в полигоне воды}\\
		\mathrm{false}, & \text{$i$ --- точка, находящаяся вне полигона воды}
	\end{cases}$\\
	Факт нахождения в полигоне воды проверяется посредством проверки нахождения точки с соответствующими координатами в полигоне воды открытой векторной геопространственной базы данных OpenStreetMap.

	\subsection*{Предварительная фильтрация}
	\textit{Данный раздел носит описательный характер. На данный момент описанные в нем фильтрации не применяются}.
	\subsubsection*{Фильтрация по количеству спутников}
	Низкое количество спутников (менее 3) в момент определения координат GPS-датчиком имеет корреляцию с $y_i = \mathrm{anomaly}$ (\textit{необходимо предоставить корреляционный анализ?}). По этой причине удаление записей $\left\{\left(\mathit{time}_i,\, \mathit{lat}_i,\, \mathit{lon}_i,\, \mathit{sat}_i\right) \right\}_{i=0,\ i \in \{sat_i < 3\}}^{n}$ уже на этапе предварительной фильтрации улучшает метрики качества.

	\subsubsection*{Фильтрация по факту нахождения в полигоне воды}
	Нахождение точки с координатами GPS-датчика вне полигона воды имеет корреляцию с $y_i = \mathrm{anomaly}$ (\textit{необходимо предоставить корреляционный анализ?}). По этой причине удаление записей \\ $\left\{\left(\mathit{time}_i,\, \mathit{lat}_i,\, \mathit{lon}_i,\, \mathit{sat}_i,\, \mathit{in\_water}_i\right) \right\}_{i=0,\ i \in \{\mathit{in\_water}_i = \mathrm{false}\}}^{n}$ уже на этапе предварительной фильтрации улучшает метрики качества.


	\section{Переход в прямоугольную локальную систему координат}
	Координаты, получаемые от GPS-датчика $(\mathit{lat}_i,\, \mathit{lon}_i)$, заданы в глобальной угловой системе WGS-84. Для моделей и метрик в плоскости координаты переводятся в локальную прямоугольную систему координат по схеме WGS-84 $\to$ ECEF $\to$ ENU.

	\subsection*{WGS-84}
	WGS-84 --- геодезическая система координат и модель Земли в виде эллипсоида вращения с большой полуосью $a=6378137$~м и сжатием $f=1/298.257223563$. Положение точки задаётся геодезическими координатами $(\varphi, \lambda, h)$: широта $\varphi \in [-90^\circ,90^\circ]$, долгота $\lambda \in [-180^\circ,180^\circ]$, высота $h$ над эллипсоидом. В телеметрии используются $\mathit{lat}_i \approx \varphi_i$, $\mathit{lon}_i \approx \lambda_i$; в эксперименте высота отсутствует, поэтому принимается $h_i = 0$.

	\subsection*{ECEF}
	Система ECEF (Earth-Centered, Earth-Fixed) --- декартова система координат в метрах с началом в центре масс Земли. Ось $Z$ направлена вдоль оси вращения эллипсоида, оси $X,Y$ лежат в экваториальной плоскости. Углы $\varphi_i, \lambda_i$ подставляются в радианах: $\varphi_i = \mathit{lat}_i \cdot \pi/180$, $\lambda_i = \mathit{lon}_i \cdot \pi/180$. Сначала вычисляется радиус кривизны первого вертикала
	\begin{equation*}
		N(\varphi_i)=\frac{a}{\sqrt{1-e^2\sin^2\varphi_i}}, \qquad e^2 = 2f - f^2.
	\end{equation*}
	Координаты точки в ECEF:
	\begin{align*}
		X_i &= \bigl(N(\varphi_i)+h_i\bigr)\cos\varphi_i\cos\lambda_i,\\
		Y_i &= \bigl(N(\varphi_i)+h_i\bigr)\cos\varphi_i\sin\lambda_i,\\
		Z_i &= \bigl(N(\varphi_i)(1-e^2)+h_i\bigr)\sin\varphi_i.
	\end{align*}
	Переход в ECEF переводит точки временного ряда в единую метрическую 3D-систему.

	\subsection*{ENU (локальная прямоугольная система)}
	Локальная система ENU (East--North--Up) задаётся относительно опорной точки $(\varphi_0,\lambda_0,h_0)$. Ось $E$ направлена на восток, ось $N$ --- на север, ось $U$ --- вверх по нормали к эллипсоиду. Разность векторов в ECEF $\Delta_i = (X_i-X_0, Y_i-Y_0, Z_i-Z_0)^{\mathsf T}$ переводится в ENU поворотом на матрицу $R(\varphi_0,\lambda_0)$:
	\begin{equation*}
		R(\varphi_0,\lambda_0)=
		\begin{pmatrix}
			-\sin\lambda_0 & \cos\lambda_0 & 0\\
			-\sin\varphi_0\cos\lambda_0 & -\sin\varphi_0\sin\lambda_0 & \cos\varphi_0\\
			\cos\varphi_0\cos\lambda_0 & \cos\varphi_0\sin\lambda_0 & \sin\varphi_0
		\end{pmatrix}, \qquad
		\begin{pmatrix}E_i\\N_i\\U_i\end{pmatrix}=R(\varphi_0,\lambda_0)\,\Delta_i.
	\end{equation*}
	Матрица $R$ ортогональна и выполняет только поворот осей без изменения масштаба. Для движения судна по поверхности используется горизонтальная проекция:
	\begin{equation*}
		o_i \equiv E_i, \qquad n_i \equiv N_i.
	\end{equation*}
	Компонента $U_i$ отбрасывается.

	\subsection*{Порядок перевода системы координат}
	\begin{enumerate}
		\item \textbf{Опорная точка.} В качестве опорной точки берётся первая валидная точка ряда после предварительной фильтрации $(\varphi_0,\lambda_0,h_0)=(\varphi_0,\lambda_0,0)$. Она переводится в ECEF: $(X_0,Y_0,Z_0)$.
		\item \textbf{Переход по точкам.} Для каждой точки $i$ вычисляются $(X_i,Y_i,Z_i)$ по формулам WGS-84$\to$ECEF, формируется разность $\Delta_i=(X_i-X_0,Y_i-Y_0,Z_i-Z_0)$ и выполняется поворот: $(E_i,N_i,U_i)^{\mathsf T}=R(\varphi_0,\lambda_0)\Delta_i$.
		\item \textbf{Проекция.} Для плоского представления траектории используются горизонтальные координаты $o_i \approx E_i$, $n_i \approx N_i$ (метры), отсчитываемые от опорной точки.
	\end{enumerate}
	Плоская интерпретация $(o_i,n_i)$ достаточно точна в окрестности опорной точки; при значительном удалении от неё начинают сказываться кривизна Земли и отличие высот.

	\subsection{Обработка временного ряда по частям}
	Обрабатываемый временной ряд может иметь большую протяженность. Обработка всего ряда целиком повышает ресурсы, используемые детектором, а также снижает точность локальной прямоугольной системы координат для точек, удаленных от опорной точки. В целях предотвращения указанных негативных факторов ряд обрабатывается по частям:
	\begin{enumerate}
		\item Временной ряд последовательно делится на интервалы по $8640$ записей (при $\Delta t = 10$~с $8640$ записей равны одним суткам). \\
		$j = 0, 1, \ldots$ --- номер интервала. В $j$-й интервал попадают индексы \\
		$\mathcal{I}_j = \{i : j \cdot 8640 \leq i \leq \min((j+1)\cdot 8640 - 1,\, n)\}$.
		\item При переходе в ENU на интервале $j$ в качестве опорной точки используется последняя точка штатного хода перед началом интервала: $j_0 = \max\{i < j \cdot 8640 : y_i = \mathrm{move}\}$. Если указанного $i$ нет, то $j_0 = 0$ (первая точка ряда).
		\item Координаты $(\varphi_0,\, \lambda_0) = (\varphi_{j_0},\, \lambda_{j_0})$ задают локальную систему для всех $i \in \mathcal{I}_j$; преобразование в $o_i$, $n_i$ выполняется в каждом интервале независимо.
	\end{enumerate}

	\section{Описание и сравнение моделей}
	\subsection{Предварительная обработка}
	При сравнении моделей исходные записи телеметрии приводятся к единому ряду наблюдений и двум сценариям разбиения на обучающую и тестовую выборки. После загрузки CSV индексы сырых строк $j=0,\ldots,M-1$; поля $t_j$ (\texttt{time}), $\varphi_j$ (\texttt{lat}), $\lambda_j$ (\texttt{lon}), $y_j$ (\texttt{status}). В эксперименте используются только $y_j\in\{\mathrm{move},\mathrm{anomaly}\}$. Для метрик классификации положительный класс --- $\mathrm{anomaly}$; числовая метка \texttt{label}: $\mathrm{anomaly}\rightarrow 0$, $\mathrm{move}\rightarrow 1$.

	Порядок действий:
	\begin{enumerate}
		\item \textbf{Отбор столбцов.} Остаются $(t_i,\varphi_i,\lambda_i,y_i)$.
		\item \textbf{Валидные координаты.}
		\[\mathcal{J}_1=\{i:\ \varphi_i,\lambda_i\ \text{конечны},\ -90\le\varphi_i\le 90,\ -180\le\lambda_i\le 180,\ (\varphi_i,\lambda_i)\neq(0,0)\}.\]
		\item \textbf{Бинарная задача.}
		\[\mathcal{J}_2=\{i\in\mathcal{J}_1:\ y_i\in\{\mathrm{move},\mathrm{anomaly}\}\}.\]
		Точки с иными метками (в том числе $\mathrm{stand}$) отбрасываются.
		\item \textbf{Упорядочивание.} На $\mathcal{J}_2$ --- сортировка по возрастанию $t_i$; индекс $i=0$ --- первая валидная точка ряда.
		\item \textbf{Локальная ENU (м).} Опорная точка --- первая валидная точка: $(\varphi_0,\lambda_0)=(\varphi_{0},\lambda_{0})$. Преобразование WGS-84 $\rightarrow$ ECEF $\rightarrow$ ENU:
		\[(o_i,n_i)=(E_i,N_i)=\mathrm{ENU}(\varphi_i,\lambda_i\mid\varphi_0,\lambda_0).\]
		Широта и долгота сохраняются; в таблицу добавляются $o_i$, $n_i$ (метка класса --- $y_i$).
		\item \textbf{Обрезка по крайним $\mathrm{move}$.}
		\[\mathcal{I}_{\mathrm{move}}=\{i:\ y_i=\mathrm{move}\},\quad i_{\min}=\min\mathcal{I}_{\mathrm{move}},\quad i_{\max}=\max\mathcal{I}_{\mathrm{move}}.\]
		Рабочий ряд --- срез $i_{\min}\le i\le i_{\max}$ (первая и последняя точки --- $\mathrm{move}$); координаты $(o_i,n_i)$ не пересчитываются.
		\item \textbf{Разбиение train / test.} Для доли обучения $\rho\in(0,1)$ задается $k_0=\lfloor\rho N\rfloor$ и граница
		\[k=\min\{i\ge k_0:\ y_i=\mathrm{move}\}.\]
		Обучение: индексы $[0,k)$, тест: $[k,N)$. Если $y_{k_0}=\mathrm{move}$, то $k=k_0$. Граница не ставится внутрь фрагмента без соседнего $\mathrm{move}$.
	\end{enumerate}

	В ходе эксперимента формируются два сценария с $\rho_1=0.6$ и $\rho_2=0.8$; для каждого представлены размеры выборок и доли $\mathrm{anomaly}$ и $\mathrm{move}$ на обучении и тесте. Разбиение по суточным интервалам с отдельной опорой ENU на каждом интервале в данном эксперименте не применяется.

	\subsection{Описание метрик}
	В сравнительных экспериментах с моделями используются коэффициент корреляции Мэттьюса (MCC) и distance loss.

	\subsubsection*{Коэффициент корреляции Мэттьюса (MCC)}
	\textbf{Коэффициент корреляции Мэттьюса} (Matthews correlation coefficient, MCC) --- коэффициент корреляции Пирсона между двумя бинарными векторами: эталонной разметкой и предсказанием (каждая компонента кодируется как 0 или 1). Эквивалентно он выражается через четыре числа таблицы сопряженности $TP$, $FP$, $FN$, $TN$ и принимает значения от $-1$ до $1$.

	MCC применяется в задачах бинарной классификации, в том числе при сильном дисбалансе классов, когда accuracy и доля «большого» класса малоинформативны. В экспериментах MCC считается по точкам; положительный класс --- аномалия ($y_i=\mathrm{anomaly}$, бинарная метка $a_i=\mathbf{1}[y_i=\mathrm{anomaly}]$; предсказание $\hat a_i=\mathbf{1}[\hat y_i=\mathrm{anomaly}]$). На обучающей выборке MCC --- главный критерий выбора порога $\tau$ (и гиперпараметров вроде $\varepsilon$, $\nu$).

	Основное преимущество: метрика одновременно учитывает все четыре ячейки таблицы сопряженности и симметрично отражает цену ложных срабатываний и пропусков; при дисбалансе классов тривиальные предсказания «всегда аномалия» или «всегда штатный ход» обычно дают MCC около нуля, тогда как accuracy может оставаться высокой.

	\begin{equation*}
		\mathrm{MCC}=\frac{TP\cdot TN - FP\cdot FN}{\sqrt{(TP+FP)(TP+FN)(TN+FP)(TN+FN)}},
	\end{equation*}
	где:\\
	$TP=\sum_i \mathbf{1}[a_i=1\land \hat a_i=1]$,\\
	$FP=\sum_i \mathbf{1}[a_i=0\land \hat a_i=1]$,\\
	$FN=\sum_i \mathbf{1}[a_i=1\land \hat a_i=0]$,\\
	$TN=\sum_i \mathbf{1}[a_i=0\land \hat a_i=0]$.

	\textbf{Интерпретация:} $\mathrm{MCC}\in[-1,1]$; $1$ --- полное совпадение масок; $0$ --- согласие на уровне случайного предсказания при тех же маргинальных долях классов; отрицательные значения --- систематическое расхождение хуже случайного. Если знаменатель равен нулю (например, в выборке или в предсказании только один класс), MCC в строгом смысле не определен; в коде экспериментов в таком случае принимается $0$.

	\textbf{Источник:} Chicco D., Jurman G. The advantages of the Matthews correlation coefficient (MCC) over F1 score and accuracy in binary classification evaluation // BMC Genomics. 2020. Vol.~21. Art.~6. Open access: \url{https://link.springer.com/article/10.1186/s12864-019-6413-7} (в статье отмечено, что MCC изначально предложен B.~W.~Matthews в 1975~г.).

	\subsubsection*{Distance loss}
	Геопространственная метрика: насколько расходятся длины эталонной и предсказанной траектории (в метрах). Длина ломаной --- сумма геодезических звеньев $d_{\mathrm{geo}}$; distance loss --- модуль разности:
	\begin{equation*}
		L(T)=\sum_i d_{\mathrm{geo}}(p_i,p_{i+1}), \qquad DL=\bigl|L_{\mathrm{pred}}-L_{\mathrm{ref}}\bigr|.
	\end{equation*}

	$L_{\mathrm{ref}}$ строится по эталону $G$ (штатные точки и линейная интерполяция в аномалиях между соседними штатными). $L_{\mathrm{pred}}$ --- по предсказанию $P$: в ломаную входят только точки, принятые за штатный ход ($\hat y_i \neq \mathrm{anomaly}$); пропуски не интерполируются, соседние принятые точки соединяются напрямую. На тесте $DL$ также считается по заполненным траекториям; при выборе порога по MCC среди кандидатов в полосе близкого значения MCC минимизируется операционный distance loss (длина по принятым точкам штатного хода).

	\subsection{Описание и тестирование моделей}
	Сравнение выполняется в блокноте на одном ряде, прошедшем вышеописанную предварительную обработку. Для каждой модели и каждого сценария ($\rho_1=0.6$, $\rho_2=0.8$): подбор гиперпараметров и порога --- только на train; на test фиксируются выбранные параметры, итоговые статистики измеряются по тестовому ряду при зафиксированных параметрах моделей. Критерий порога по score(функции правдоподобия): максимум MCC на train; среди кандидатов в полосе $0{,}02$ от лучшего MCC --- минимум distance loss (длина ломаной по принятым точкам $\mathrm{move}$ в ENU). Итоговые таблицы на test содержат метрики из раздела~«Список метрик» (без MCC и train-метрик). Исключения: DBSCAN (порог --- радиус $\varepsilon$, маска шума), One-Class SVM (порог score $s>0$, подбор $\nu$), MPCT (сетка $v_+$ и $\tau_{\mathrm{shape}}$, бинарное правило без $\tau$ по score).

	\paragraph{Выбор порога}
	На обучающей выборке: сетка кандидатов $\tau$ --- квантили оценки аномальности и 60 значений, равномерных в логарифмической шкале между её 1-м перцентилем и максимумом; выбирается максимум MCC по точкам с оценкой, среди кандидатов в полосе MCC (до 0.02 от лучшего) --- наименьшая операционная ошибка длины. На тестовой выборке --- тот же порог без подбора.

	К сетке добавляется канонический кандидат, если у оценки аномальности есть теоретическая граница: $\chi^2_{0.999}(2)$ у моделей с квадратом нормированной инновации (фильтр Калмана, гауссовский процесс, GP-EKF), 3 у фильтра Хэмпела, 0.5 у Isolation Forest и HMM, 0 у One-Class SVM. Собственная сетка вместо порога --- у DBSCAN (радиус $\varepsilon$) и MPCT (предел скорости $v_+$); у LOF порог подбирается вместе с числом соседей $k$.

	Отвергнутая точка не используется в оценках следующих точек. У моделей прогноза (AR, ARIMA, фильтр Калмана, гауссовский процесс, GP-EKF) она не входит в прогноз; у моделей по признакам (фильтр Хэмпела, DBSCAN, LOF, Isolation Forest, One-Class SVM, Robust PCA, LSTM-автоэнкодер, TCN, DONUT, HMM) решение принимается в два прохода: признаки пересчитываются от последних принятых точек, оценки и решения второго прохода --- итоговые. Отвергнутые точки зависят от порога, поэтому замена выполняется заново для каждого кандидата. MPCT замены не требует: скорость считается только между выбранными точками.

	\paragraph{Точки без оценки и покрытие}
	Точка без оценки --- точка, в которой модель не может дать оценку аномальности; она не относится ни к штатному ходу, ни к аномалии. Точность, полнота, F1 и MCC считаются только по точкам с оценкой, у каждой модели --- на своём подмножестве. В заполненной траектории точка без оценки восстанавливается интерполяцией по времени, как отвергнутая. Задержка обнаружения отсчитывается от начала эпизода аномалии. Покрытие в итоговых таблицах --- доля точек без оценки, доля аномальных точек среди них и доля аномальных точек без оценки.

	\subsubsection{MPCT (Maximum Physically Consistent Trajectory)}

	\paragraph{Постановка}

	В каждом отрезке ряда выбирается самая длинная подпоследовательность точек, в которой судно могло пройти от каждой выбранной точки к следующей со скоростью не выше $v_+$. Аномалия --- точка, не вошедшая в эту подпоследовательность. Модель пакетная (batch): выбор точки зависит и от следующих точек.

	\paragraph{Модель и решение}

	Подпоследовательность $i_1<i_2<\ldots<i_m$ точек отрезка допустима, если для соседних выбранных точек
	\[
	t_{i_{k+1}}-t_{i_k}>180\ \text{с}\quad\text{или}\quad \frac{\lVert\mathbf p_{i_{k+1}}-\mathbf p_{i_k}\rVert}{t_{i_{k+1}}-t_{i_k}}\le v_+,\qquad \mathbf p_i=(o_i,n_i).
	\]

	Выбирается допустимая подпоследовательность наибольшей длины; $\hat a_i=1$ для точек вне неё. Расстояние --- в плоскости ENU, как у остальных моделей. Через 180~с без выбранных точек цепочка начинается заново --- как перезапуск у остальных моделей.

	Решение находится динамическим программированием: для каждой точки --- длина лучшей допустимой подпоследовательности, оканчивающейся в ней; предшественник ищется среди точек за последние 180~с или, при перезапуске, среди более ранних. При равной длине выбирается ближайший по времени предшественник.

	Глобальный выбор устраняет зависимость от одной опорной точки: аномальная точка не может «увести» решение, потому что цепочка через неё короче. Выброс не оставляет следа на соседних точках --- скорость считается только между выбранными точками. Решение бинарное, непрерывной оценки аномальности нет; роль порога играет $v_+$.

	Свойство модели: смещённый эпизод длиннее 180~с, физически согласованный внутри, может оказаться выбран вместо штатных точек, если он длиннее их.

	\paragraph{Время и точки без оценки}

	Отрезки --- между разрывами по времени больше 180~с. Из точек с одинаковым временем остаётся первая, остальные --- без оценки; все остальные точки, включая первую точку отрезка, получают решение.

	\paragraph{Обучение и выбор порога}

	$v_+$ --- из 30 значений, равномерных в логарифмической шкале от 95-го процентиля скоростей между соседними точками штатного хода обучающей выборки до трёх их максимумов: ниже 95-го процентиля отвергался бы обычный ход, выше трёх максимумов модель не отвергает ничего. Выбор --- как порог у AR(2).

	\subsubsection{AR(2)}

	\paragraph{Постановка}
	Авторегрессия второго порядка описывает уровень ряда --- плоскую координату --- через два предыдущих значения. Аномалия --- точка с большой ошибкой причинного одношагового прогноза.

	\paragraph{Модель и оценка аномальности}
	Для восточной координаты (аналогично для северной):
	\[
	o_j=\mu_o+\eta_j,\qquad \eta_j=\phi_1\eta_{j-1}+\phi_2\eta_{j-2}+\varepsilon_j,\qquad \varepsilon_j\sim\mathcal N(0,\sigma_o^2),
	\]
	\[
	\hat o_j=\mu_o+\phi_1(o_{j-1}-\mu_o)+\phi_2(o_{j-2}-\mu_o),
	\]
	где $\mu_o$ --- первая точка отрезка.

	Оценка аномальности --- ошибка прогноза, м:
	\[
	s_j=\sqrt{(o_j-\hat o_j)^2+(n_j-\hat n_j)^2},\qquad \hat a_j=\mathbf{1}[s_j>\tau].
	\]

	Отвергнутая точка в лаги не входит: следующий прогноз строится от последних принятых точек.

	\paragraph{Время и точки без оценки}
	Ряд переносится на регулярную сетку по времени с шагом 10~с, $j$ --- узел сетки. В узлах без точек (пропусках) прогноз продолжается без обновления, поэтому прогноз точки после пропуска строится на $h$ шагов от последней принятой. Прогноз начинается заново после разрыва по времени больше 180~с и после 180~с подряд отвергнутых точек; первые две точки нового отрезка --- без оценки. Оценка аномальности не нормируется на $h$ и после пропусков и серий отвергнутых точек систематически выше.

	\paragraph{Обучение и выбор порога}
	\begin{enumerate}
		\item Параметры $\mu$, $\phi_1$, $\phi_2$, $\sigma^2$ каждой оси оцениваются по максимуму правдоподобия только по точкам штатного хода обучающей выборки: аномальные точки --- пропуски, правдоподобие суммируется по отрезкам между сериями пропусков длиннее 180~с.
		\item Оценка аномальности на обучающей выборке считается по всем точкам.
		\item Сетка кандидатов порога строится по оценке аномальности без замены отвергнутых точек; для каждого кандидата детектор заново проходит обучающую выборку с заменой.
	\end{enumerate}

	\subsubsection{ARIMA(2,1,1)}
	Как у AR(2).

	\paragraph{Постановка}
	Приращения координаты описываются процессом авторегрессии и скользящего среднего; уровень --- их накопление.

	\paragraph{Модель и оценка аномальности}
	Для восточной координаты (аналогично для северной):
	\[
	\Delta o_j=\phi_1\Delta o_{j-1}+\phi_2\Delta o_{j-2}+\varepsilon_j+\theta_1\varepsilon_{j-1},\qquad \varepsilon_j\sim\mathcal N(0,\sigma_o^2),
	\]
	где $\Delta o_j=o_j-o_{j-1}$. Дрейф не вводится: средняя скорость по обучающей выборке физического смысла не имеет. Член $\theta_1\varepsilon_{j-1}$ описывает шум GPS: ошибка измерения входит в два соседних приращения с разными знаками.

	\paragraph{Время и точки без оценки}
	Первые три точки отрезка --- без оценки: одна на приращение, две на лаги.

	\subsubsection{Фильтр Калмана (постоянная скорость)}

	\paragraph{Постановка}
	Скрытое состояние --- положение и скорость в плоскости ENU, измерение --- положение. Аномалия --- точка, плохо согласованная с причинным прогнозом фильтра.

	\paragraph{Модель и оценка аномальности}
	\[
	\mathbf x_i=(o_i,n_i,v^o_i,v^n_i)^\top,\qquad \mathbf x_i=F_i\mathbf x_{i-1}+\mathbf w_i,\qquad \mathbf z_i=H\mathbf x_i+\boldsymbol\nu_i,
	\]
	\[
	F_i=\begin{pmatrix}I&\Delta t_i I\\0&I\end{pmatrix},\qquad Q_i=\sigma_{\mathrm{acc}}^2\begin{pmatrix}\tfrac{\Delta t_i^4}{4}I&\tfrac{\Delta t_i^3}{2}I\\\tfrac{\Delta t_i^3}{2}I&\Delta t_i^2 I\end{pmatrix},\qquad R=\sigma_{\mathrm{meas}}^2 I,
	\]
	где $I$ --- единичная матрица $2\times2$, $H=(I\;0)$, $\Delta t_i=t_i-t_{i-1}$.

	Оценка аномальности --- квадрат нормированной инновации:
	\[
	s_i=\mathbf r_i^\top S_i^{-1}\mathbf r_i,\qquad \mathbf r_i=\mathbf z_i-H\hat{\mathbf x}_{i\mid i-1},\qquad S_i=HP_{i\mid i-1}H^\top+R,\qquad \hat a_i=\mathbf{1}[s_i>\tau].
	\]

	Отвергнутая точка состояние не обновляет: следующий прогноз строится от последней принятой точки. Аномальная точка, принятая после серии отвергнутых, искажает оценку скорости: следующие точки отвергаются до перезапуска прогноза.

	\paragraph{Время и точки без оценки}
	Как у AR(2): прогноз начинается заново после разрыва по времени больше 180~с и после 180~с подряд отвергнутых точек; первые две точки отрезка --- без оценки; из точек с одинаковым временем остаётся первая. Начальное состояние --- по первым двум точкам отрезка:
	\[
	\hat{\mathbf x}=\Bigl(\mathbf z_2,\ \frac{\mathbf z_2-\mathbf z_1}{\Delta t}\Bigr),\qquad P=\begin{pmatrix}R&R/\Delta t\\R/\Delta t&2R/\Delta t^2\end{pmatrix}.
	\]

	\paragraph{Обучение и выбор порога}
	\begin{enumerate}
		\item $\sigma_{\mathrm{acc}}$ и $\sigma_{\mathrm{meas}}$ оцениваются по максимуму правдоподобия только по точкам штатного хода обучающей выборки: аномальные точки пропускаются, правдоподобие суммируется по отрезкам между сериями пропусков длиннее 180~с, первые две точки отрезка в него не входят. Нижние границы --- $\sigma_{\mathrm{acc}}\ge10^{-3}$~м/с$^2$, $\sigma_{\mathrm{meas}}\ge0.5$~м: при меньших значениях правдоподобие численно вырождается.
		\item Оценка аномальности и подбор порога --- как у AR(2); в сетку кандидатов добавляется $\chi^2_{0.999}(2)$.
	\end{enumerate}

	\subsubsection{Гауссовский процесс}

	\paragraph{Постановка}
	Гауссовский процесс по окну последних принятых точек прогнозирует положение следующей точки. Аномалия --- точка, плохо согласованная с прогнозом.

	\paragraph{Модель и оценка аномальности}
	Для восточной координаты (аналогично для северной) на окне из последних 18 принятых точек отрезка:
	\[
	o(t)=\beta_0+\beta_1(t-\bar t)+f(t),\qquad f\sim\mathcal{GP}\bigl(0,k(t,t')\bigr),\qquad k(t,t')=\sigma_{\mathrm{rbf}}^2\exp\Bigl(-\frac{(t-t')^2}{2\ell^2}\Bigr)+\sigma_{\mathrm n}^2\,\delta_{tt'},
	\]
	где $\bar t$ --- среднее время точек окна. Линейный тренд описывает движение с постоянной скоростью; $\beta_0$, $\beta_1$ оцениваются обобщённым МНК по окну без априорного распределения. RBF описывает плавный манёвр, $\sigma_{\mathrm n}^2$ --- шум GPS.

	Прогноз в момент $t_i$ --- апостериорные среднее $\mu_i$ и дисперсия $\sigma_i^2$, включающая неопределённость оценки тренда. Оценка аномальности --- квадрат нормированной ошибки по двум осям:
	\[
	s_i=\sum_{\text{оси}}\frac{(z_i-\mu_i)^2}{\sigma_i^2+\sigma_{\mathrm n}^2},\qquad \hat a_i=\mathbf{1}[s_i>\tau].
	\]

	Отвергнутая точка в окно не входит. Свойства модели: аномальная точка, принятая в окно, искажает тренд, пока не выйдет из окна; после паузы, близкой к разрыву по времени, прогноз продолжает тренд окна без учёта изменения скорости за паузу. В обоих случаях следующие точки могут отвергаться до перезапуска прогноза.

	\paragraph{Время и точки без оценки}
	Как у AR(2): прогноз начинается заново после разрыва по времени больше 180~с и после 180~с подряд отвергнутых точек; из точек с одинаковым временем остаётся первая. Первые три точки отрезка --- без оценки; пока в окне меньше 18 точек, прогноз строится по имеющимся.

	\paragraph{Обучение и выбор порога}
	\begin{enumerate}
		\item Гиперпараметры $\sigma_{\mathrm{rbf}}$, $\ell$, $\sigma_{\mathrm n}$ --- общие для обеих осей; оцениваются по максимуму суммы логарифмов прогнозной плотности следующей точки на окнах штатного хода обучающей выборки.
		\item Оценка аномальности и подбор порога --- как у AR(2); в сетку кандидатов добавляется $\chi^2_{0.999}(2)$.
	\end{enumerate}

	\subsubsection{GP-EKF (фильтр Калмана с гауссовским процессом)}
	Как у фильтра Калмана (постоянная скорость); ниже --- только отличия.

	\paragraph{Постановка}
	Модель постоянной скорости дополняется ускорением-поправкой, которая зависит от скорости и обучается гауссовским процессом по штатному ходу. Аномалия --- точка, плохо согласованная с причинным прогнозом фильтра.

	\paragraph{Модель и оценка аномальности}
	\[
	\mathbf x_i=F_i\mathbf x_{i-1}+B_i\,\mathbf g(\mathbf v_{i-1})+\mathbf w_i,\qquad B_i=\begin{pmatrix}\tfrac{\Delta t_i^2}{2}I\\\Delta t_i I\end{pmatrix},\qquad \mathbf v_i=(v^o_i,v^n_i)^\top,
	\]
	где $\mathbf g=(g^o,g^n)^\top$ --- ускорение-поправка, м/с$^2$. По каждой оси --- независимый гауссовский процесс от скорости; для восточной оси (аналогично для северной):
	\[
	g^o\sim\mathcal{GP}\bigl(0,k^o(\mathbf v,\mathbf v')\bigr),\qquad k^o(\mathbf v,\mathbf v')=\sigma_{f,o}^2\exp\Bigl(-\frac12\sum_{m\in\{o,n\}}\frac{(v^m-v'^m)^2}{\ell_{o,m}^2}\Bigr),
	\]
	где $\sigma_{f,o}$ --- амплитуда, $\ell_{o,m}$ --- своя длина корреляции по каждой составляющей скорости.

	Прогноз --- расширенный фильтр Калмана (extended Kalman filter, EKF):
	\[
	\hat{\mathbf x}_{i\mid i-1}=F_i\hat{\mathbf x}_{i-1}+B_i\boldsymbol\mu_i,\qquad G_i=F_i+B_i\,(0\;\;J_i),\qquad P_{i\mid i-1}=G_iP_{i-1}G_i^\top+B_i\Sigma_iB_i^\top,
	\]
	где $\boldsymbol\mu_i$ и $\sigma_{g,o}^2$, $\sigma_{g,n}^2$ --- апостериорные среднее и дисперсии гауссовских процессов при скорости $\hat{\mathbf v}_{i-1}$, $J_i=\partial\boldsymbol\mu/\partial\mathbf v$ --- якобиан среднего, $\Sigma_i=\mathrm{diag}\bigl(\sigma_{g,o}^2+\sigma_{\mathrm{acc},o}^2,\ \sigma_{g,n}^2+\sigma_{\mathrm{acc},n}^2\bigr)$, $\sigma_{\mathrm{acc},o}$, $\sigma_{\mathrm{acc},n}$ --- шум ускорения.

	Шум процесса растёт с $\Delta t_i$, как у фильтра Калмана, и больше при скоростях, где мало опорных точек. При $\mathbf g\equiv0$ и $\sigma_{\mathrm{acc},o}=\sigma_{\mathrm{acc},n}$ модель совпадает с фильтром Калмана; вне диапазона скоростей опорных точек $\boldsymbol\mu_i\to0$.

	Свойство модели: сглаживатель притягивает обучающие переходы к модели постоянной скорости, поэтому поправка занижена; невоспроизведённая часть ускорения входит в шум ускорения $\sigma_{\mathrm{acc},o}$, $\sigma_{\mathrm{acc},n}$.

	Оценка аномальности, решение и обновление состояния --- как у фильтра Калмана.

	\paragraph{Обучение и выбор порога}
	\begin{enumerate}
		\item Опорные точки. Сглаживатель Рауха --- Тунга --- Штрибеля (RTS smoother) с моделью и оценками $\sigma_{\mathrm{acc}}$, $\sigma_{\mathrm{meas}}$ фильтра Калмана восстанавливает состояния $\tilde{\mathbf x}_j$ на отрезках штатного хода обучающей выборки. Опорная точка --- переход: вход $\tilde{\mathbf v}_{j-1}$, цель $(\tilde{\mathbf v}_j-\tilde{\mathbf v}_{j-1})/\Delta t_j$. Цель по сырым разностям координат определялась бы шумом GPS. Опорные точки разнесены по времени не меньше чем на 120~с, из них берётся до 1000 случайных: ошибки сглаживания соседних переходов коррелированы и принимались бы за зависимость от скорости.
		\item $\sigma_f$ и $\ell$ --- по максимуму маргинального правдоподобия (marginal likelihood) на опорных точках; шум цели оценивается вместе с ними и в $\Sigma_i$ не входит.
		\item $\sigma_{\mathrm{acc},o}$, $\sigma_{\mathrm{acc},n}$, $\sigma_{\mathrm{meas}}$ --- по максимуму правдоподобия фильтра при обученных гауссовских процессах, как $\sigma_{\mathrm{acc}}$ и $\sigma_{\mathrm{meas}}$ у фильтра Калмана, с теми же нижними границами. Сглаживатель убирает часть шума процесса, поэтому оценка шума ускорения по опорным точкам занижена.
		\item Оценка аномальности и подбор порога --- как у фильтра Калмана.
	\end{enumerate}

	При прогоне выводится доля дисперсии ускорения штатного хода, объяснённая средним $\boldsymbol\mu$, на обучающей и тестовой выборках --- мера отличия модели от фильтра Калмана.

	\subsubsection{Скрытая марковская модель (HMM)}

	Признаки, нормировка и второй проход --- как у DBSCAN; ниже --- только отличия.

	\paragraph{Постановка}

	Ряд признаков описывается одной скрытой марковской цепью, в которой есть режимы штатного хода и режимы аномалии. Переходы между режимами, в том числе между классами, выучиваются по разметке обучающей выборки. Аномалия --- точка, в которой вероятнее режим аномалии.

	\paragraph{Модель и оценка аномальности}
	\[
	S_i\in\mathcal S_{\mathrm{move}}\cup\mathcal S_{\mathrm{anom}},\qquad |\mathcal S_{\mathrm{move}}|=|\mathcal S_{\mathrm{anom}}|=3,\qquad P(S_1=k)=\pi_k,\qquad P(S_i=k\mid S_{i-1}=j)=T_{jk},\qquad \mathbf u_i\mid S_i=k\sim\mathcal N(\mathbf m_k,D_k),
	\]
	где $\mathcal S_{\mathrm{move}}$, $\mathcal S_{\mathrm{anom}}$ --- режимы штатного хода и аномалии, $\mathbf m_k$ и $D_k$ --- среднее и диагональная ковариация признаков в режиме $k$.

	Оценка аномальности --- причинная апостериорная вероятность режима аномалии:
	\[
	f_i(k)\propto\mathcal N(\mathbf u_i;\mathbf m_k,D_k)\sum_j f_{i-1}(j)\,T_{jk},\qquad s_i=\sum_{k\in\mathcal S_{\mathrm{anom}}}f_i(k),\qquad \hat a_i=\mathbf{1}[s_i>\tau],
	\]
	где $f_i(k)=P(S_i=k\mid\mathbf u_{1:i})$, в начале отрезка $f(k)\propto\pi_k\,\mathcal N(\mathbf u;\mathbf m_k,D_k)$.

	Вероятности перехода между классами задают типичную длительность эпизодов: внутри длинного эпизода аномалии цепь остаётся в режимах аномалии, даже если признаки отдельной точки похожи на штатный ход. В отличие от остальных моделей, HMM учитывает длительность эпизодов и использует разметку при обучении. Свойство модели: после одиночного выброса цепь некоторое время остаётся в режимах аномалии, и несколько следующих точек тоже размечаются аномалией.

	\paragraph{Время и точки без оценки}

	Фильтр начинается заново в каждом отрезке между разрывами по времени больше 180~с. Без оценки: первые две точки отрезка (как у DBSCAN); из точек с одинаковым временем остаётся первая.

	\paragraph{Обучение и выбор порога}

	\begin{enumerate}
		\item Нормировка --- как у DBSCAN. Признаки при обучении считаются по последней точке штатного хода, чтобы след аномальной точки не попадал в признаки соседней.
		\item $\pi$, $T$, $\mathbf m_k$, $D_k$ --- алгоритмом EM (Баума --- Уэлча) по всем отрезкам обучающей выборки; в размеченной точке вероятность режимов чужого класса равна нулю. Итерации --- до относительного прироста логарифма правдоподобия меньше $10^{-4}$, не больше 200; из 5 запусков с фиксированными начальными значениями генератора выбирается лучший по правдоподобию. Число режимов не подбирается.
		\item Подбор порога --- как у AR(2); в сетку кандидатов добавляется 0.5 --- граница «аномалия вероятнее штатного хода»; второй проход выполняется для каждого кандидата.
	\end{enumerate}

	\subsubsection{Фильтр Хэмпела}

	\paragraph{Постановка}
	По модулю скорости в скользящем окне оцениваются медиана и типичный разброс. Аномалия --- точка, скорость которой далеко от медианы окна. Модель пакетная (batch): окно включает следующие точки.

	\paragraph{Модель и оценка аномальности}
	\[
	v_i=\frac{\sqrt{\Delta o_i^2+\Delta n_i^2}}{\Delta t_i},\qquad m_i=\mathrm{med}(W_i),\qquad \mathrm{MAD}_i=\mathrm{med}_{v\in W_i}|v-m_i|,
	\]
	\[
	s_i=\frac{|v_i-m_i|}{\max\bigl(1.4826\,\mathrm{MAD}_i,\ \sigma_{\min}\bigr)},\qquad \hat a_i=\mathbf{1}[s_i>\tau],
	\]
	где $W_i$ --- скорости $2w+1$ точек отрезка вокруг точки $i$, $w=10$; у краёв отрезка --- ближайшие $2w+1$ точек, в коротком отрезке --- весь отрезок; $1.4826\,\mathrm{MAD}_i$ --- оценка СКО при нормальном распределении; $\sigma_{\min}$ --- нижняя граница разброса.

	Медиана устойчива, пока аномальных точек в окне не больше $w$. Свойство модели: эпизод аномалии длиннее $w$ точек сдвигает медиану и не обнаруживается; при постоянном смещении скорость аномальна только на входе в эпизод и выходе из него.

	Одиночный выброс искажает и скорость следующей точки. Поэтому решение принимается в два прохода: после первого скорости пересчитываются от последней принятой точки (отвергнутые точки пропускаются), оценки и решения второго прохода --- итоговые.

	\paragraph{Время и точки без оценки}
	Окно не выходит за отрезок между разрывами по времени больше 180~с. Без оценки: первая точка отрезка (скорость не определена); из точек с одинаковым временем остаётся первая; все точки отрезка, в котором меньше пяти скоростей.

	\paragraph{Обучение и выбор порога}
	\begin{enumerate}
		\item $\sigma_{\min}$ --- 5-й перцентиль $1.4826\,\mathrm{MAD}_i$ по точкам штатного хода обучающей выборки. Нижняя граница нужна, когда скорости в окне почти одинаковы (стоянка, округление координат): иначе малое отклонение даёт огромную оценку аномальности.
		\item $w$ не подбирается. Подбор порога --- как у AR(2); в сетку кандидатов добавляется 3, для каждого кандидата второй проход выполняется заново.
	\end{enumerate}

	\subsubsection{DBSCAN}

	\paragraph{Постановка}
	Точка описывается кинематическими признаками; плотные области признаков штатного хода обучающей выборки задают норму. Аномалия --- точка, далёкая от плотной области.

	\paragraph{Модель и оценка аномальности}
	\[
	\mathbf f_i=\bigl(v_i,\ a^{\mathrm{long}}_i,\ a^{\mathrm{lat}}_i\bigr),\qquad a^{\mathrm{long}}_i=\frac{v_i-v_{i-1}}{\Delta t_i},\qquad a^{\mathrm{lat}}_i=\frac{v_i\,\Delta\psi_i}{\Delta t_i},\qquad \mathbf u_i=\frac{\mathbf f_i-m}{\mathrm{IQR}},
	\]
	где $v_i$ --- скорость на шаге от предыдущей точки, $\Delta\psi_i$ --- изменение курса между соседними шагами, $m$ и $\mathrm{IQR}$ --- медиана и межквартильный размах признаков на штатном ходу обучающей выборки. Признаки не зависят от курса и от $\Delta t$; поперечное ускорение мало на стоянке, где курс случаен.

	Ядровая точка (core point) --- точка штатного хода обучающей выборки, в $\varepsilon$-окрестности которой не меньше $\mathrm{MinPts}=6$ точек штатного хода, включая её саму. Оценка аномальности --- расстояние до ближайшей ядровой точки:
	\[
	s_i=\min_{c\in C_\varepsilon}\lVert\mathbf u_i-\mathbf u_c\rVert,\qquad \hat a_i=\mathbf{1}[s_i>\varepsilon],
	\]
	где $C_\varepsilon$ --- множество ядровых точек. Порог совпадает с радиусом: точка в $\varepsilon$-окрестности ядровой точки относится к её кластеру (правило DBSCAN). Решение по точке не зависит от остальных точек выборки.

	Выброс искажает признаки следующих точек. Поэтому решение принимается в два прохода: после первого признаки пересчитываются от последних принятых точек (отвергнутые точки пропускаются), оценки и решения второго прохода --- итоговые.

	\paragraph{Время и точки без оценки}
	Признаки считаются внутри отрезка между разрывами по времени больше 180~с. Без оценки: первые две точки отрезка (для ускорения нужны две скорости); из точек с одинаковым временем остаётся первая.

	\paragraph{Обучение и выбор порога}
	\begin{enumerate}
		\item $m$, $\mathrm{IQR}$ и ядровые точки --- по точкам штатного хода обучающей выборки: признаки считаются по соседним точкам штатного хода, аномальные точки пропускаются, отрезки --- между сериями пропусков длиннее 180~с.
		\item $\varepsilon$ --- из 30 квантилей расстояния до 6-го соседа на штатном ходу обучающей выборки (уровни 0.5--0.999). Уровень $q$ --- доля ядровых точек среди точек штатного хода: при $q<0.5$ область нормы сжимается до самого плотного ядра и отвергается заметная часть штатного хода; при $q=1$ кандидат задаёт самая изолированная точка штатного хода (например, ошибка разметки). Для каждого кандидата ядровые точки находятся заново, разметка обучающей выборки --- в два прохода. Выбор --- как порог у AR(2). $\mathrm{MinPts}$ не подбирается.
	\end{enumerate}

	\subsubsection{Local Outlier Factor (LOF)}
	Признаки, нормировка, точки без оценки и второй проход --- как у DBSCAN; ниже --- только отличия.

	\paragraph{Постановка}
	Локальная плотность точки в пространстве признаков сравнивается с плотностью её соседей среди точек штатного хода обучающей выборки (обнаружение новизны, novelty detection). Аномалия --- точка в области, разреженной относительно соседей.

	\paragraph{Модель и оценка аномальности}
	Для набора $T$ нормированных признаков точек штатного хода обучающей выборки:
	\[
	\mathrm{rd}_k(\mathbf u,\mathbf y)=\max\bigl(d_k(\mathbf y),\ \lVert\mathbf u-\mathbf y\rVert\bigr),\qquad \mathrm{lrd}_k(\mathbf u)=\Bigl(\frac1k\sum_{\mathbf y\in N_k(\mathbf u)}\mathrm{rd}_k(\mathbf u,\mathbf y)\Bigr)^{-1},
	\]
	\[
	s_i=\mathrm{LOF}_k(\mathbf u_i)=\frac1k\sum_{\mathbf y\in N_k(\mathbf u_i)}\frac{\mathrm{lrd}_k(\mathbf y)}{\mathrm{lrd}_k(\mathbf u_i)},\qquad \hat a_i=\mathbf{1}[s_i>\tau],
	\]
	где $N_k(\mathbf u)$ --- $k$ ближайших к $\mathbf u$ точек набора $T$, $d_k(\mathbf y)$ --- расстояние от $\mathbf y$ до её $k$-го соседа в $T$, $\mathrm{rd}_k$ --- расстояние достижимости (reachability distance), $\mathrm{lrd}_k$ --- локальная плотность достижимости (local reachability density). $s_i\approx1$ --- плотность как у соседей; $s_i\gg1$ --- точка в разреженной области. Для точки из $T$ сама точка в $N_k$ не входит.

	В отличие от DBSCAN, оценка аномальности непрерывна и относительна: плотная область «стоянка» и разреженная область «ход» оцениваются каждая относительно своих соседей. Решение по точке не зависит от остальных точек выборки.

	Свойство модели (риск принят): точки с совпадающими признаками (стоянка с повторяющимися координатами) дают нулевые расстояния достижимости и очень большую $\mathrm{lrd}_k$; у соседней точки с обычным шумом GPS $s_i$ тогда велика --- возможны ложные срабатывания рядом со стоянками.

	\paragraph{Обучение и выбор порога}
	\begin{enumerate}
		\item Нормировка и набор $T$ --- как у DBSCAN.
		\item $k$ и порог подбираются совместно: $k\in\{10,20,35,50\}$; для каждого $k$ --- общая сетка кандидатов порога по оценке первого прохода, второй проход на каждого кандидата; выбор по всем парам --- как порог у AR(2). Свойство модели: оценка аномальности штатного хода слабо зависит от $k$, поэтому выбор $k$ определяется малыми различиями MCC и может меняться между сценариями.
	\end{enumerate}

	\subsubsection{Isolation Forest}
	Признаки, нормировка, точки без оценки и второй проход --- как у DBSCAN; ниже --- только отличия.

	\paragraph{Постановка}
	Лес случайных деревьев, обученный на штатном ходу обучающей выборки, изолирует точку последовательными разрезами по признакам (обнаружение новизны). Аномалия --- точка, которая отделяется за малое число разрезов.

	\paragraph{Модель и оценка аномальности}
	Каждое из 200 деревьев строится по случайной подвыборке из $n_{\mathrm{sub}}=256$ точек набора $T$ (как у LOF): в узле выбираются случайный признак и случайное значение разреза в диапазоне точек узла.
	\[
	s_i=2^{-\bar h(\mathbf u_i)/c(n_{\mathrm{sub}})},\qquad c(n_{\mathrm{sub}})=2H(n_{\mathrm{sub}}-1)-\frac{2(n_{\mathrm{sub}}-1)}{n_{\mathrm{sub}}},\qquad \hat a_i=\mathbf{1}[s_i>\tau],
	\]
	где $\bar h(\mathbf u_i)$ --- средняя по деревьям длина пути от корня до листа точки (с поправкой $c$ на число точек в листе), $H(n)\approx\ln n+0.5772$ --- гармоническое число, $c(n_{\mathrm{sub}})$ --- средняя длина пути в двоичном дереве поиска из $n_{\mathrm{sub}}$ точек. Точка вне области штатного хода уходит в крайнюю ветку и изолируется за несколько разрезов: $s_i$ близка к 1; у точки штатного хода $s_i<0.5$.

	В отличие от DBSCAN и LOF, осевые разрезы оценивают положение точки относительно всей области штатного хода, а не локальную плотность. Решение по точке не зависит от остальных точек выборки.

	\paragraph{Обучение и выбор порога}
	\begin{enumerate}
		\item Нормировка и набор $T$ --- как у DBSCAN; лес обучается на $T$ при фиксированном начальном значении генератора; число деревьев и $n_{\mathrm{sub}}$ не подбираются.
		\item Подбор порога --- как у AR(2); в сетку кандидатов добавляется 0.5 --- граница аномальности из определения оценки; второй проход выполняется для каждого кандидата.
	\end{enumerate}

	\subsubsection{One-Class SVM}
	Признаки, нормировка, точки без оценки и второй проход --- как у DBSCAN; ниже --- только отличия.

	\paragraph{Постановка}
	По точкам штатного хода обучающей выборки строится граница области штатного хода в пространстве признаков (обнаружение новизны). Аномалия --- точка, далёкая от этой области.

	\paragraph{Модель и оценка аномальности}
	\[
	f(\mathbf u)=\sum_j\alpha_j K(\mathbf u_j,\mathbf u)-\rho,\qquad K(\mathbf u,\mathbf u')=\exp\bigl(-\gamma\lVert\mathbf u-\mathbf u'\rVert^2\bigr),\qquad \gamma=\frac{1}{d\,\mathrm{Var}(T)},
	\]
	\[
	s_i=-f(\mathbf u_i),\qquad \hat a_i=\mathbf{1}[s_i>\tau],
	\]
	где сумма --- по опорным векторам $\mathbf u_j$ набора $T$ (как у LOF), $\alpha_j\ge0$ и $\rho$ --- решение задачи One-Class SVM при $\nu=0.1$, $d=3$ --- число признаков, $\mathrm{Var}(T)$ --- дисперсия всех значений признаков набора $T$. Внутри области штатного хода $f>0$; при $\tau=0$ вне области оказывается не больше доли $\nu$ точек $T$.

	$\nu$ ограничивает снизу долю опорных векторов: при $\nu=0.1$ их не меньше 10\% точек $T$, и оценка ближе к сглаженной оценке плотности --- информативна не только у границы, но и во всём пространстве признаков.

	В отличие от DBSCAN, LOF и Isolation Forest, граница гладкая и определяется взвешенной суммой ядер опорных векторов. Решение по точке не зависит от остальных точек выборки.

	Свойство модели: далеко от области штатного хода значения ядра близки к нулю и оценка насыщается на $\rho$ --- далёкие точки между собой не различаются; на решение это не влияет, пока порог ниже $\rho$.

	\paragraph{Обучение и выбор порога}
	\begin{enumerate}
		\item Нормировка и набор $T$ --- как у DBSCAN; модель обучается на всех точках $T$; $\nu$ и $\gamma$ не подбираются.
		\item Подбор порога --- как у AR(2); в сетку кандидатов добавляется 0 --- граница SVM; второй проход выполняется для каждого кандидата.
	\end{enumerate}

	\subsubsection{Robust PCA (Principal Component Pursuit)}

	Признаки, нормировка и второй проход --- как у DBSCAN; ниже --- только отличия.

	\paragraph{Постановка}

	Окно из подряд идущих точек раскладывается на часть, объяснимую типичным поведением штатного хода, и разреженное отклонение. Аномалия --- точка, на которую приходится отклонение.

	\paragraph{Модель и оценка аномальности}

	Окно $W$ --- вектор $\mathbf x_W=(\mathbf u_j,\ldots,\mathbf u_{j+n_W-1})$ нормированных признаков $n_W=16$ подряд идущих точек.

	Типичное поведение --- подпространство окон штатного хода обучающей выборки. Матрица $X$ этих окон (по строке на окно) раскладывается на низкоранговую и разреженную части:
	\[
	\min_{L,S}\ \lVert L\rVert_*+\lambda\lVert S\rVert_1\quad\text{при}\quad X=L+S,\qquad \lambda=\frac{1}{\sqrt{\max(n,m)}},
	\]
	где $\lVert L\rVert_*$ --- сумма сингулярных чисел (ядерная норма), $\lVert S\rVert_1$ --- сумма модулей элементов, $n\times m$ --- размер $X$. Базис подпространства $U$ --- правые сингулярные векторы $L$ с ненулевыми сингулярными числами; ранг определяется разложением. Единичные выбросы и ошибки разметки штатного хода уходят в $S$ и подпространство не искажают.

	Оценка окна любой выборки --- разреженный остаток проекции на фиксированное подпространство:
	\[
	\mathbf c^*=\arg\min_{\mathbf c}\lVert\mathbf x_W-U\mathbf c\rVert_1,\qquad \mathbf s_W=\mathbf x_W-U\mathbf c^*,\qquad s_i=\max_{W\ni i}\lVert\mathbf s_W(i)\rVert,\qquad \hat a_i=\mathbf{1}[s_i>\tau],
	\]
	где $\mathbf s_W(i)$ --- три компоненты остатка, относящиеся к точке $i$; максимум --- по окнам, содержащим точку. Это разложение $X=L+S$ при известном подпространстве: остаток разрежен без отдельного штрафа. Отклонение одной точки не объясняется комбинацией окон штатного хода и попадает в остаток на позициях этой точки; остаток обычной проекции размазал бы его по всему окну.

	В отличие от точечных моделей, оценивается согласованность точки с соседями: аномальна последовательность, даже если значения признаков по отдельности обычны. Свойство модели: эпизод длиннее окна с обычной кинематикой не отличается от штатного хода; аномальны окна, захватывающие вход в эпизод и выход из него.

	\paragraph{Время и точки без оценки}

	Окна --- по $n_W$ подряд идущих точек внутри отрезка между разрывами по времени больше 180~с, с шагом 1. Без оценки: первые две точки отрезка (как у DBSCAN); все точки отрезка, в котором меньше $n_W+2$ точек.

	\paragraph{Обучение и выбор порога}

	\begin{enumerate}
		\item Нормировка и окна штатного хода --- по точкам штатного хода обучающей выборки, как у DBSCAN; разложение $X=L+S$ --- до сходимости.
		\item $n_W$ не подбирается. Подбор порога --- как у AR(2); второй проход выполняется для каждого кандидата.
	\end{enumerate}

	\subsubsection{LSTM-автоэнкодер (EncDec-AD)}

	Признаки, нормировка и второй проход --- как у DBSCAN; окна и точки без оценки --- как у Robust PCA; ниже --- только отличия.

	\paragraph{Постановка}

	Сеть кодировщик–декодер на LSTM, обученная на окнах штатного хода обучающей выборки, восстанавливает окно признаков. Аномалия --- точка, признаки которой восстанавливаются плохо.

	\paragraph{Модель и оценка аномальности}

	Кодировщик LSTM читает окно $\mathbf x_W=(\mathbf u_j,\ldots,\mathbf u_{j+n_W-1})$ и сжимает его в скрытое состояние размера 32; декодер LSTM восстанавливает из этого состояния окно $\hat{\mathbf x}_W$. Для точки $i$ на позиции $p$ окна $W$:
	\[
	\mathbf e_W(i)=\mathbf u_i-\hat{\mathbf u}_W(i),\qquad d_W(i)=\sum_{k=1}^{3}\frac{\bigl(e_{W,k}(i)-\mu_{p,k}\bigr)^2}{\sigma_{p,k}^2},\qquad s_i=\min_{W\ni i}d_W(i),\qquad \hat a_i=\mathbf{1}[s_i>\tau],
	\]
	где $\hat{\mathbf u}_W(i)$ --- восстановленные признаки точки $i$, $\mu_{p,k}$ и $\sigma_{p,k}^2$ --- среднее и дисперсия ошибки признака $k$ на позиции $p$ по калибровочным окнам; минимум --- по окнам, содержащим точку.

	Сеть восстанавливает окно из сжатого состояния, поэтому выброс портит восстановление и соседних точек того же окна. Выброс восстанавливается плохо во всех своих окнах, а у соседней точки есть окно без выброса --- минимум по окнам относит отклонение к точке выброса.

	В отличие от Robust PCA, норма --- нелинейная модель последовательности, а не линейное подпространство. Свойство модели --- как у Robust PCA: эпизод длиннее окна с обычной кинематикой не отличается от штатного хода.

	\paragraph{Обучение и выбор порога}

	\begin{enumerate}
		\item Нормировка --- как у DBSCAN. Окна штатного хода обучающей выборки упорядочиваются по времени: последние 15\% --- калибровочные, окна, пересекающиеся с ними, отбрасываются; из остальных берётся до 8000 случайных окон для обучения.
		\item Сеть обучается по среднему квадрату ошибки восстановления (Adam), до 50 эпох с ранней остановкой: 5 эпох без улучшения ошибки на калибровочных окнах; начальное значение генератора фиксировано, вычисления детерминированы. Размер сети и число эпох не подбираются.
		\item $\mu_{p,k}$, $\sigma_{p,k}^2$ --- по калибровочным окнам. Подбор порога --- как у AR(2); второй проход выполняется для каждого кандидата.
	\end{enumerate}

	\subsubsection{Temporal Convolutional Network (TCN)}
	Признаки, нормировка и второй проход --- как у DBSCAN; обучение и калибровка --- как у LSTM-автоэнкодера; ниже --- только отличия.

	\paragraph{Постановка}
	Причинная свёрточная сеть по признакам предыдущих принятых точек прогнозирует признаки следующей точки. Аномалия --- точка, признаки которой далеки от прогноза.

	\paragraph{Модель и оценка аномальности}
	\[
	\hat{\mathbf u}_i=g\bigl(\mathbf u_{i-15},\ldots,\mathbf u_{i-1}\bigr),\qquad s_i=\sum_{k=1}^{3}\frac{\bigl(u_{i,k}-\hat u_{i,k}-\mu_k\bigr)^2}{\sigma_k^2},\qquad \hat a_i=\mathbf{1}[s_i>\tau],
	\]
	где $g$ --- сеть из трёх причинных свёрточных уровней (ядро 3, расширения 1, 2, 4, 32 канала), её поле зрения --- ровно 15 предыдущих точек; $\mu_k$, $\sigma_k^2$ --- среднее и дисперсия ошибки прогноза признака $k$ по калибровочным окнам. Шаг по времени до прогнозируемой точки входит в её признаки: скорость и ускорения нормированы по $\Delta t$.

	В истории --- только принятые точки: история --- 15 последних принятых точек отрезка, отвергнутая точка в неё не входит.

	В отличие от LSTM-автоэнкодера, модель прогнозирует следующую точку по прошлому, а не восстанавливает окно; в отличие от AR и фильтра Калмана --- нелинейная модель по признакам, а не линейная модель положения.

	\paragraph{Время и точки без оценки}
	Без оценки: первые 17 точек отрезка (2 --- на признаки, 15 --- на историю); из точек с одинаковым временем остаётся первая.

	\paragraph{Обучение и выбор порога}
	\begin{enumerate}
		\item Окно обучения --- 15 точек истории и прогнозируемая точка штатного хода; калибровка, число окон, обучение и ранняя остановка --- как у LSTM-автоэнкодера.
		\item $\mu_k$, $\sigma_k^2$ --- по калибровочным окнам. Подбор порога --- как у AR(2); второй проход выполняется для каждого кандидата.
	\end{enumerate}

	\subsubsection{DONUT (вариационный автоэнкодер)}

	Признаки, нормировка и второй проход --- как у DBSCAN; окно и точки без оценки --- как у TCN; ниже --- только отличия.

	\paragraph{Постановка}

	Вариационный автоэнкодер задаёт распределение окна признаков штатного хода. Обучение --- на всех окнах обучающей выборки, аномальные точки исключаются из правдоподобия (modified ELBO, M-ELBO). Аномалия --- точка с малой вероятностью восстановления.

	\paragraph{Модель и оценка аномальности}

	Окно $\mathbf x_W=(\mathbf u_{i-15},\ldots,\mathbf u_i)$ --- 15 предыдущих принятых точек и точка $i$:
	\[
	z\sim\mathcal N(0,I),\qquad \mathbf u_j\mid z\sim\mathcal N\bigl(\boldsymbol\mu_j(z),\ \mathrm{diag}\,\boldsymbol\sigma_j^2(z)\bigr),\qquad q(z\mid\mathbf x_W)=\mathcal N\bigl(\boldsymbol\mu_q(\mathbf x_W),\ \mathrm{diag}\,\boldsymbol\sigma_q^2(\mathbf x_W)\bigr),
	\]
	где $z$ --- скрытый вектор размерности 8; декодер ($\boldsymbol\mu_j$, $\boldsymbol\sigma_j$) и кодировщик ($\boldsymbol\mu_q$, $\boldsymbol\sigma_q$) --- по два полносвязных слоя из 100 нейронов; $\boldsymbol\sigma_j\ge10^{-4}$.

	Обучение --- максимум M-ELBO:
	\[
	\mathcal L(\mathbf x_W)=\mathbb E_{q(z\mid\mathbf x_W)}\Bigl[\sum_{j\in W}\alpha_j\log p(\mathbf u_j\mid z)\Bigr]-\beta\,D_{\mathrm{KL}}\bigl(q(z\mid\mathbf x_W)\,\Vert\,\mathcal N(0,I)\bigr),\qquad \beta=\frac{1}{|W|}\sum_{j\in W}\alpha_j,
	\]
	где $\alpha_j=1$ у точки штатного хода и $\alpha_j=0$ у аномальной; при обучении ещё 1\% точек случайно получает $\alpha_j=0$ --- модель учится восстанавливать окно с невидимыми точками. Разметка обучающей выборки используется уже при обучении, а не только при выборе порога.

	Оценка аномальности --- минус вероятность восстановления последней точки окна:
	\[
	s_i=-\frac{1}{L}\sum_{l=1}^{L}\log p(\mathbf u_i\mid z_l),\qquad z_l\sim q(z\mid\mathbf x_W),\qquad L=100,\qquad \hat a_i=\mathbf{1}[s_i>\tau].
	\]

	Во втором проходе отвергнутые точки в окно не входят, поэтому восстанавливать пропущенные точки при оценке не требуется.

	В отличие от LSTM-автоэнкодера и TCN, модель вероятностная: оценка --- логарифм правдоподобия, отдельная калибровка ошибки не нужна; разметка участвует в обучении.

	\paragraph{Обучение и выбор порога}

	\begin{enumerate}
		\item Окна --- все окна обучающей выборки, кроме окон из одних аномальных точек; признаки при обучении считаются по последней точке штатного хода, чтобы след аномальной точки не попадал в признаки соседней. Калибровочный блок, число окон, ранняя остановка (по M-ELBO) и начальное значение генератора --- как у LSTM-автоэнкодера.
		\item Подбор порога --- как у AR(2); второй проход выполняется для каждого кандидата.
	\end{enumerate}

	\subsubsection{STGVAD --- исключён из сравнения}
	STGVAD строит граф состояний нескольких судов: пространственные рёбра связывают близкие суда, аномалия --- поведение судна, необычное относительно соседей. В данных одно судно без идентификатора, поэтому граф взаимодействия судов построить нельзя. Граф по окну одной траектории вырождается: все прошлые точки окна связаны между собой, и модель сводится к модели последовательности, как TCN и LSTM-автоэнкодер. Модель исключена из сравнения; вернуть её можно при данных нескольких судов в одном окне времени.

\subsection{Итоговые таблицы на тесте (сортировка)}
Сводка по моделям блокнота на тестовых выборках сценариев $\rho_1=0{,}6$ и $\rho_2=0{,}8$. Ниже --- таблицы, отсортированные по убыванию $F_1$ и по возрастанию distance loss (DL).
\subsubsection{Сортировка: F1}
\begin{table}[H]
\centering
\scriptsize
\setlength{\tabcolsep}{3pt}
\caption{Сортировка: F1: df1 — доля обучения $\rho=0{,}6$ (тестовая выборка)}
\begin{tabular}{lrrrrrrrrrrrrr}
\toprule
Модель & RMSE, м & Max(d\_i), м & Distance loss, м & Delay median, с & Delay mean, с & Delay p95, с & Delay max, с & On-time episodes & Late/missed episodes & Precision & Recall & F1 & MCC \\
\midrule
MPCT & & & & & & & & & & & & & \\
One-Class SVM & & & & & & & & & & & & & \\
Isolation Forest & & & & & & & & & & & & & \\
Фильтр Калмана & & & & & & & & & & & & & \\
HMM & & & & & & & & & & & & & \\
DONUT & & & & & & & & & & & & & \\
LSTM-автоэнкодер & & & & & & & & & & & & & \\
Robust PCA & & & & & & & & & & & & & \\
GP-EKF & & & & & & & & & & & & & \\
TCN & & & & & & & & & & & & & \\
Gaussian process & & & & & & & & & & & & & \\
ARIMA(2,1,1) & & & & & & & & & & & & & \\
AR(p) & & & & & & & & & & & & & \\
LOF & & & & & & & & & & & & & \\
Hampel & & & & & & & & & & & & & \\
\bottomrule
\end{tabular}
\end{table}

\begin{table}[H]
\centering
\scriptsize
\setlength{\tabcolsep}{3pt}
\caption{Сортировка: F1: df2 — доля обучения $\rho=0{,}8$ (тестовая выборка)}
\begin{tabular}{lrrrrrrrrrrrrr}
\toprule
Модель & RMSE, м & Max(d\_i), м & Distance loss, м & Delay median, с & Delay mean, с & Delay p95, с & Delay max, с & On-time episodes & Late/missed episodes & Precision & Recall & F1 & MCC \\
\midrule
MPCT & & & & & & & & & & & & & \\
Isolation Forest & & & & & & & & & & & & & \\
One-Class SVM & & & & & & & & & & & & & \\
Фильтр Калмана & & & & & & & & & & & & & \\
HMM & & & & & & & & & & & & & \\
Robust PCA & & & & & & & & & & & & & \\
DONUT & & & & & & & & & & & & & \\
GP-EKF & & & & & & & & & & & & & \\
LSTM-автоэнкодер & & & & & & & & & & & & & \\
TCN & & & & & & & & & & & & & \\
AR(p) & & & & & & & & & & & & & \\
ARIMA(2,1,1) & & & & & & & & & & & & & \\
LOF & & & & & & & & & & & & & \\
Hampel & & & & & & & & & & & & & \\
Gaussian process & & & & & & & & & & & & & \\
\bottomrule
\end{tabular}
\end{table}

\subsubsection{Сортировка: Distance loss}
\begin{table}[H]
\centering
\scriptsize
\setlength{\tabcolsep}{3pt}
\caption{Сортировка: Distance loss: df1 — доля обучения $\rho=0{,}6$ (тестовая выборка)}
\begin{tabular}{lrrrrrrrrrrrrr}
\toprule
Модель & RMSE, м & Max(d\_i), м & Distance loss, м & Delay median, с & Delay mean, с & Delay p95, с & Delay max, с & On-time episodes & Late/missed episodes & Precision & Recall & F1 & MCC \\
\midrule
MPCT & & & & & & & & & & & & & \\
Фильтр Калмана & & & & & & & & & & & & & \\
Isolation Forest & & & & & & & & & & & & & \\
GP-EKF & & & & & & & & & & & & & \\
AR(p) & & & & & & & & & & & & & \\
One-Class SVM & & & & & & & & & & & & & \\
HMM & & & & & & & & & & & & & \\
ARIMA(2,1,1) & & & & & & & & & & & & & \\
Gaussian process & & & & & & & & & & & & & \\
LOF & & & & & & & & & & & & & \\
Hampel & & & & & & & & & & & & & \\
LSTM-автоэнкодер & & & & & & & & & & & & & \\
DONUT & & & & & & & & & & & & & \\
TCN & & & & & & & & & & & & & \\
Robust PCA & & & & & & & & & & & & & \\
\bottomrule
\end{tabular}
\end{table}

\begin{table}[H]
\centering
\scriptsize
\setlength{\tabcolsep}{3pt}
\caption{Сортировка: Distance loss: df2 — доля обучения $\rho=0{,}8$ (тестовая выборка)}
\begin{tabular}{lrrrrrrrrrrrrr}
\toprule
Модель & RMSE, м & Max(d\_i), м & Distance loss, м & Delay median, с & Delay mean, с & Delay p95, с & Delay max, с & On-time episodes & Late/missed episodes & Precision & Recall & F1 & MCC \\
\midrule
MPCT & & & & & & & & & & & & & \\
Isolation Forest & & & & & & & & & & & & & \\
Фильтр Калмана & & & & & & & & & & & & & \\
One-Class SVM & & & & & & & & & & & & & \\
GP-EKF & & & & & & & & & & & & & \\
HMM & & & & & & & & & & & & & \\
AR(p) & & & & & & & & & & & & & \\
ARIMA(2,1,1) & & & & & & & & & & & & & \\
LOF & & & & & & & & & & & & & \\
LSTM-автоэнкодер & & & & & & & & & & & & & \\
Hampel & & & & & & & & & & & & & \\
DONUT & & & & & & & & & & & & & \\
Gaussian process & & & & & & & & & & & & & \\
TCN & & & & & & & & & & & & & \\
Robust PCA & & & & & & & & & & & & & \\
\bottomrule
\end{tabular}
\end{table}

\subsection{Разбор результатов, ревью и выводы}

Числа ниже сняты с сохраненного прогона общей сводки (16 строк). В каталоге блокнота семнадцать детекторов; строки DBSCAN в этом снимке нет, поэтому плотности кластеров в выводах нет. Два сценария — префиксы одного очищенного ряда, а не независимые суда.

\subsubsection{1. Устройство данных}

Исходный CSV: 419\,453 строк. После отбора координат, статусов move и anomaly, сортировки по времени и обрезки по крайним точкам move остаётся 252\,818 наблюдений.

\begin{table}[H]
\centering
\caption{Размеры выборок и доли классов по сценариям}
\begin{tabular}{lrrrrr}
\toprule
Сценарий & $\rho$ & Обучение & Anomaly, \% (train) & Тест & Anomaly, \% (test) \\
\midrule
df1 & & & & & \\
df2 & & & & & \\
\bottomrule
\end{tabular}
\end{table}

Тест df2 — хвост того же ряда, что и тест df1. Согласие рангов на двух $\rho$ говорит об устойчивости вдоль трека.

\subsubsection{2. Как устроены модели}

Общая схема: по ряду строится score $s_i$ (или сразу бинарная маска), затем $\hat a_i=\mathbf{1}[s_i>\tau]$. Признаки почти нигде не включают абсолютные координаты карты. Исключения по смыслу — положение в ENU у линейных моделей и фильтров и положение относительно начала окна у нейросетей. Типичный кинематический вектор: шаг на восток и север, скорость, ускорение, приращение курса, иногда логарифм формы окна длины 6.

По тому, что модель считает «нормой», методы распадаются на шесть групп.

\textbf{Физическое ограничение.} MPCT не учит распределение. Якорь — последняя принятая точка move. Следующая точка принимается, если скорость от якоря не выше $v_+$ и форма окна не выше порога. Отклонённая точка якорь не сдвигает, поэтому дальнее облако само не становится новым штатным участком. Это жадная адаптация идеи longest physically consistent subsequence, а не глобальный поиск подпоследовательности из статьи.

\textbf{Линейный прогноз координат.} AR($p$) с $p=2$ и линейным трендом и ARIMA(2, 1, 1) оцениваются по всему обучению, включая anomaly, отдельно по востоку и северу. Score — евклидова ошибка одношагового прогноза. Порядок фиксирован, а не выбирается по AIC на скользящем окне, как в исходной схеме Qin et al. для ARIMA.

\textbf{Фильтр состояния.} Фильтр Калмана ведёт положение и скорость в модели постоянной скорости. Шум процесса — дискретное белошумное ускорение, шаг в матрице перехода ограничен 45 с. Score — normalized innovation squared. При разрыве больше 180 с состояние сбрасывается на измерение. GP-EKF ставит на место приращения постоянной скорости четыре независимых GP по переходам move$\to$move (скорость и $\Delta t$, не координаты на карте). Наблюдение остаётся линейным, $\sigma_z=40$ м. Оба фильтра причинные.

\textbf{Гауссовский процесс на шаге.} Два независимых GP описывают $(\Delta x,\Delta y)$ как функцию времени по подвыборке точек move. Score — сумма нормированных квадратов остатков. Вдали от опор прогноз сходится к среднему шагу, поэтому смена района сама по себе score не поднимает.

\textbf{Модель обоих классов.} HMM — три гауссовских HMM с диагональной ковариацией на робастно нормированных $(v, a^{\mathrm{acc}}, \Delta\psi)$: штатные отрезки, аномальные отрезки и быстрые аномальные отрезки ($v\ge 8$ м/с). Score — разность причинных предсказательных лог-правдоподобий anomaly и move. Это единственный детектор, которому на обучении нужны оба статуса как отдельные режимы. Остальные либо одноклассовые, либо смотрят на смесь, либо вообще не оценивают параметры по выборке.

\textbf{Локальные и признаковые детекторы без динамики судна.} Hampel сравнивает скорость с медианой центрированного окна ($h=10$): метод рассчитан на одиночные выбросы и видит будущее внутри окна. LOF и Isolation Forest работают с робастными кинематическими признаками; лес строится по всему обучению и применяется к тесту, LOF пересчитывается на размечаемом ряду. One-Class SVM с ядром RBF учится только на подвыборке move (не больше 4000 точек), аномалия — точка вне оценённого носителя. Robust PCA раскладывает матрицу окон длины 16 на низкий ранг и разреженную часть; score точки — максимум нормы разреженной компоненты по покрывающим окнам.

\textbf{Оконные нейросети на штатных окнах.} LSTM-автоэнкодер восстанавливает окно положения в километрах от первой точки окна; score — диагональное расстояние Махаланобиса ошибки, точке пишется максимум по окнам. TCN причинно предсказывает следующую точку того же окна. DONUT — VAE тех же окон, обучение через Pyro SVI; в score входят ошибка реконструкции в среднем энкодера и KL. STGVAD на окне одной последовательности (в статье граф нескольких судов) предсказывает положение по скорости, курсу и рёбрам короче 5 км. Общее свойство: окно, а не шаг фильтра, и агрегация max по окнам, которые накрывают точку.

В блокноте описаны, но не прогнаны bootstrap particle filter (мультимодальный аналог Калмана) и Anomaly Transformer (ошибка восстановления, умноженная на расхождение attention и prior-ассоциации).

\subsubsection{3. Тестирование}

Ранги по $F_1$ (убывание) и по distance loss (возрастание). DL указан в миллионах метров, чтобы порядки читались рядом.

\begin{table}[H]
\centering
\caption{Сводка рангов по $F_1$ и distance loss (выдержка)}
\footnotesize
\begin{tabular}{clrrclrr}
\toprule
Место & df1, $F_1$ & $F_1$ & DL, млн м & df2, $F_1$ & $F_1$ & DL, млн м \\
\midrule
1 & & & & & & \\
2 & & & & & & \\
3 & & & & & & \\
4 & & & & & & \\
5 & & & & & & \\
6 & & & & & & \\
16 & & & & & & \\
\bottomrule
\end{tabular}
\end{table}

По distance loss первая тройка на обоих сценариях одна и та же: MPCT, фильтр Калмана, Isolation Forest. На df1 следом GP-EKF (22.3 млн м) при $F_1$ только 0.593. На df2 четвёртое-пятое места делят One-Class SVM и GP-EKF (около 27 млн м).

Isolation Forest даёт лучший RMSE на обоих тестах: 44.3 км (df1) и 39.2 км (df2). У MPCT RMSE 73.2 и 92.5 км при лучшем DL. Калман — 63.6 и 79.0 км. Линейные модели на df2 выходят на RMSE 652 км при том же Max$(d_i)$, что и у нейросетей.

Задержка отделяет «попал в метку» от «успел на эпизод». HMM на df1 срабатывает почти сразу: медиана 10 с, 95-й перцентиль 19 с, 26 on-time против 3 late/missed. Isolation Forest стабилен: медиана ровно 110 с на обоих сценариях, late/missed — 1 эпизод из 29 и из 13. MPCT точнее по $F_1$, но опаздывает чаще: 11 и 5 late/missed, хвост задержки до 720 с. У Калмана медиана терпимая (65 и 120 с), зато среднее около 1000–1300 с и 95-й перцентиль в тысячи секунд: часть эпизодов находится очень поздно. LSTM и DONUT опаздывают на часы (медианы от 770 с до 3170 с, максимум 33\,552 с).

Сдвиг $\rho$ с 60\% на 80\% почти не двигает лидеров и ломает два метода. $F_1$ MPCT 0.845 $\to$ 0.833, Isolation Forest 0.777 $\to$ 0.833, Калман 0.756 $\to$ 0.731, HMM 0.712 $\to$ 0.705. Gaussian process 0.503 $\to$ 0.094 при precision 0.985 и recall 0.050: детектор становится почти глухим. AR($p$) и ARIMA на df2 совпадают до последней цифры (precision 0.310, recall 0.9997, $F_1$ 0.473, задержки не определены). Доля anomaly на этом тесте — 31.0\%. Совпадение precision с этой долей при полном recall означает, что обе модели пометили аномалией практически каждую точку теста. На df1 они ещё различаются ($F_1$ 0.440 и 0.416) и обе имеют задержку в часы.

\subsubsection{4. Сильные и слабые стороны по группам}

\textbf{MPCT.} Лучший $F_1$ и лучший DL на обоих разрезах, precision и recall одновременно выше 0.82. Метод прямо кодирует то, чем размеченный anomaly, судя по результату, и отличается от move: скачок скорости относительно последнего доверенного положения и изломанная форма короткого окна. Ему не мешает, что на обучении anomaly — большинство: оцениваются два порога, а не плотность смеси. Слабые места — эвристический якорь вместо глобально оптимальной подпоследовательности, фиксированная геометрия окна и более слабый RMSE, чем у леса. Часть эпизодов он пропускает или берет с задержкой до 12 минут. Имеет смысл как основной физический детектор или как ограничитель для любой статистической маски: точка, которую MPCT отвергает, не должна молча оставаться в ломаной.

\textbf{Isolation Forest.} Второй устойчивый результат и лучшая геометрия точек (RMSE, на df2 ещё и Max$(d_i)$ 92 км против сотен у остальных). Recall 0.90 и 0.93, precision 0.68 и 0.75: лес скорее помечает лишнее, чем пропускает. Задержка короткая и одинаковая на обоих хвостах. Лес видит смесь классов на обучении и всё равно отделяет редкий кинематический режим на тесте, где anomaly уже не большинство. Ограничение: нет модели движения и нет причинности внутри дерева; признак формы и шага не отличает «судно повернуло штатно, но редко» от поломки трека. DL хуже Калмана на df1 (11.4 против 8.1 млн м) при лучшем $F_1$, то есть лишние ложные anomaly режут ломаную не так болезненно, как пропуски, но и не дотягивают до MPCT.

\textbf{Фильтр Калмана.} Лучший компромисс среди моделей с состоянием: $F_1$ около 0.73–0.76, DL сразу за MPCT, precision на df2 равен 0.87 при recall 0.63. Постоянная скорость хорошо описывает куски реального хода и плохо — долгий аномальный режим, пока инновация не накопится: отсюда тяжёлый хвост задержки. Сброс на паузе 180 с честно обрабатывает разрыв связи, но обнуление скорости после паузы само может дать всплеск score. Улучшение здесь предметное: несколько режимов движения (IMM), отдельный режим «стоп», ограничение на принятие точки по правилу MPCT.

\textbf{One-Class SVM.} Высокий $F_1$ за счёт баланса precision и recall, на df1 задержка умеренная (медиана 60 с, 24 on-time из 29). На df2 $F_1$ проседает до 0.75, Max$(d_i)$ около 1000 км, DL около 27 млн м. Граница RBF на подвыборке 4000 точек move чувствительна к тому, какие штатные манёвры попали в опору. Фиксированный порог $f<0$ не участвует в полосе MCC, поэтому операционный DL при выборе $\nu$ и длины ядра может расходиться с тем, как выбирался $\tau$ у соседей.

\textbf{HMM.} Единственная модель с явной гипотезой «аномалия — другой марковский режим». Задержка на df1 лучшая в опыте, $F_1$ стабильно около 0.71. Precision 0.65 и 0.75 при recall 0.78 и 0.66: путаница симметричнее, чем у нейросетей. Геометрия слабая (RMSE 174 и 230 км, DL 36 и 33 млн м): метка эпизода ставится вовремя, но принятые и отвергнутые точки всё ещё плохо собирают длину.

\textbf{GP и GP-EKF.} Процесс на шаге слишком осторожен: precision 0.99 на обоих тестах, recall 0.34 и затем 0.05. Длинная корреляция по времени и обучение только на move дают узкую трубку типичного шага; всё, что в неё попало, принимается, остальное почти не встречается на тесте df2. GP-EKF честнее как фильтр: $F_1$ около 0.59 на обоих разрезах, DL лучше голого GP (22 и 27 млн м), recall всё ещё около 0.5. Четыре независимых скалярных GP, не больше 160 опор и сетка гиперпараметров по правдоподобию не компенсируют отсутствие связи координат и бедность опоры.

\textbf{AR и ARIMA.} На всём обучении, где anomaly составляет 60–70\%, коэффициенты подстраиваются под сломанные шаги. На длинном тесте оба прогноза перестают отличать режимы и сваливаются в одну маску «всё anomaly». Одинаковые цифры на df2 — симптом вырождения порога, а не совпадения двух теорий. Задержки в часы на df1 показывают, что даже там одношаговая ошибка уровня координат реагирует слишком поздно. Линейный прогноз уровня ENU для этой разметки слабый.

\textbf{Hampel и LOF.} $F_1$ 0.27 и 0.26 у Hampel, 0.37 и 0.44 у LOF, recall Hampel около 0.19. Локальная плотность и медиана скорости ловят изолированные скачки. Размеченная anomaly на этом ряду — протяжённый режим: внутри режима точка похожа на соседей, и LOF/Hampel её отпускают. Задержка у пойманных эпизодов при этом короткая (у LOF медиана 45 и 10 с).

\textbf{Robust PCA.} На df2 $F_1$ 0.67 выше всех оконных сетей, recall 0.87, но precision 0.54 и DL 69 млн м. Разреженная компонента окон не совпадает с геометрией принятой ломаной: низкий ранг забирает плавный ход, а в $S$ остаётся и штатный манёвр, и кусок anomaly. Max$(d_i)$ снова на общем потолке тысяч километров.

\textbf{LSTM, TCN, DONUT, STGVAD.} Общий портрет: precision высокий (часто 0.87–0.96), recall 0.33–0.58, DL на порядок хуже Калмана, задержка эпизода — тысячи секунд, кроме более быстрого STGVAD (медиана 45 и 40 с) ценой худшего DL в опыте. Окно относительно собственной первой точки хорошо видит «это окно не похоже на штатные окна» и плохо покрывает длинный аномальный участок: max по окнам зажигает точку только если какое-то окно целиком странное, а короткие окна внутри долгого плохого режима могут выглядеть внутренне согласованными. Обучение только на move при доле anomaly на обучении около двух третей ещё и оставляет мало штатных окон.

\subsubsection{5. Ревью постановки}

Опыт отвечает на узкий вопрос: на одном очищенном треке, при временном разрезе и при пороге по MCC на обучении, какой детектор ближе к статусной разметке anomaly и какой меньше портит длину принятой ломаной. На этот вопрос ответ устойчивый.

Переносить его шире мешают следующие свойства протокола.

Обучение и тест — один трек. Доля anomaly на обучении вдвое выше, чем на тесте. Любой метод, который оценивает «норму» по всей обучающей выборке (AR, ARIMA, Isolation Forest, отчасти DBSCAN и LOF на своём ряду), видит anomaly как частый режим. То, что лес всё равно работает, усиливает результат. То, что линейные модели ломаются, ожидаемо при такой смеси.

Метрика выбора порога и метрики отчёта не совпадают. $\tau$ максимизирует MCC на обучении. В таблицах теста MCC нет, зато есть $F_1$ и DL. Часть разрыва между $F_1$ One-Class SVM и DL Калмана может быть ценой этого выбора, а не только ценой семейства.

Batch и центрированные окна (Hampel, LOF по тестовому облаку, RPCA, max по окнам LSTM и DONUT, DBSCAN по тестовому облаку) используют будущее относительно точки $i$. Сравнивать их задержку с Калманом, HMM, TCN и MPCT как «скорость онлайн-детектора» нельзя без пометки: у batch-метода задержка — свойство маски, а не задержка решения в потоке.

Адаптации расходятся со статьями: фиксированный ARIMA без скользящего AIC, GP на шаге, а не на положении, HMM с непрерывными наблюдениями вместо пяти дискретных моделей Toloue и Jahan, STGVAD на одном судне, DONUT без modified ELBO для пропусков. Отрицательный результат дискредитирует конкретную реализацию в блокноте. Он не закрывает исходную статью.

Нет второго судна, нет блока по типу anomaly (скачок, залипание, подмена координат, длинный дрейф), нет интервала на разнице $F_1$. Два $\rho$ коррелированы. Совпадение Max$(d_i)$ у пяти моделей — повод проверить, не один ли и тот же непомеченный участок тянет хорду, а не пять разных механизмов.

DBSCAN в сводку этого прогона не вошёл. Пока ячейка не собрана в ту же таблицу, вывод «плотностные методы слабы» опирается только на LOF.

\subsubsection{6. Что имеет смысл рассматривать дальше}

Имеет смысл развивать, в таком порядке.

1. MPCT как несущий слой: глобальная physically consistent subsequence вместо жадного якоря; отдельный порог на длинную паузу; связка «статистический score обязан согласиться с якорем», прежде чем точка попадёт в ломаную.
2. Isolation Forest и близкие дешёвые модели на том же векторе признаков (случайный лес и градиентный бустинг, если разметка обучения доступна; иначе лес остаётся сильным unsupervised-ориентиром). Имеет смысл разобрать вклад shape, ускорения и курса, а не наращивать глубину сети.
3. Фильтр Калмана с несколькими режимами (IMM: ход, поворот, стоп) и сбросом, согласованным с MPCT. Сюда же bootstrap particle filter из methods.md: он нужен, если у судна есть развилки курса, которые один гауссовый прогноз размазывает. На текущих цифрах это следующий фильтр после IMM, не вместо Калмана.
4. HMM, только в сценарии, где оба статуса размечены. Имеет смысл больше состояний и явная длительность режима. Без разметки anomaly этот пункт закрывается.
5. One-Class SVM оставить как лёгкий признаковый контроль. Подбор $\nu$ вести той же полосой MCC и операционного DL, что и $\tau$ у остальных, иначе сравнение порогов неравное.

Имеет смысл проверить точечно, не строя вокруг этого новую линейку.

\begin{itemize}
\item GP-EKF с мультивыходным остатком поверх постоянной скорости и с опорой заметно больше 160 переходов. Голый GP по времени шага повторять не стоит: recall на df2 уже показывает отказ.
\item Авторегрессию приращений только на отрезках move, с переоценкой на коротком окне. Текущие AR и ARIMA на уровнях по всей выборке повторять не стоит.
\item Правила на разрывы $\Delta t$ отдельно от «аномальной кинематики», чтобы задержка и DL не смешивали дыру связи и порчу координат.
\end{itemize}
Не выглядит оправданным в ближайший заход.

\begin{itemize}
\item Ещё один оконный автоэнкодер того же вида (LSTM, VAE, TCN на относительном положении с max-агрегацией). Потолок precision при провале recall и общий Max$(d_i)$ уже воспроизведены четыре раза.
\item STGVAD, пока нет нескольких судов в одном окне. На одном треке граф не добавляет того, что обещает статья.
\item Hampel и LOF как основные детекторы режима. Их место — фильтр одиночных сбоев до входа в Калман или MPCT.
\item Anomaly Transformer, пока оконный score не переведён на точку и на причинный префикс. Иначе он попадёт в ту же яму, что LSTM и DONUT, только дороже.
\end{itemize}
Отдельный следующий эксперимент важнее новой архитектуры: разрезать anomaly на типы (изолированный выброс, длинный режим, разрыв времени) и посчитать $F_1$ и DL по типам. Тогда станет видно, не выигрывает ли MPCT один тип, а HMM — другой. Второй обязательный контроль — хотя бы ещё один трек или судно. Пока его нет, фраза «модель лучше» означает «на хвостах этого ряда».

\subsubsection{7. Выводы}

MPCT лучше всех совпадает с разметкой и меньше всех портит длину ломаной. Isolation Forest лучше всех держит RMSE и почти не пропускает эпизоды, с устойчивой задержкой около двух минут. Фильтр Калмана — лучший онлайн-фильтр с состоянием: чуть слабее по $F_1$, заметно сильнее по DL, чем признаковые и оконные альтернативы, с тяжёлым хвостом задержки. HMM уместен, если есть оба класса, и даёт самую быструю реакцию ценой геометрии. One-Class SVM силён как классификатор признаков и слаб как источник принятой траектории.

Линейный прогноз уровней по смеси классов на длинном тесте вырождается в метку «всё anomaly». Гауссовский процесс шага вырождается в обратную сторону: почти ничего не объявляет аномалией. Локальные выбросовые детекторы не видят протяжённый режим. Оконные нейросети в текущей адаптации устойчиво дают высокую precision, низкий recall, многочасовую задержку и хорды в тысячи километров.


\end{document}
